% =============================================================================
% Relazione progetto OOP 2025/26 -- Aurea Mediocritas
% Gruppo: Dapporto Pietro, Piraino Nikita, Conventi Mario, Locatelli Federico
% Struttura conforme alla metarelazione (Pianini, 14 aprile 2026)
%
% Compilazione consigliata:
%   pdflatex report.tex
%   pdflatex report.tex   (seconda passata per indice e riferimenti)
% =============================================================================

\documentclass[11pt, a4paper, oneside]{report}

% --- Pacchetti di base -------------------------------------------------------
\usepackage[utf8]{inputenc}
\usepackage[T1]{fontenc}
\usepackage[italian]{babel}
\usepackage{lmodern}
\usepackage[a4paper, margin=2.5cm]{geometry} 
\usepackage{microtype} 
\usepackage{amsmath}

% --- Grafica, link, codice ---------------------------------------------------
\usepackage{graphicx}
\graphicspath{{img/}}
\usepackage{float}
\usepackage{hyperref}
\usepackage{amssymb}
\hypersetup{
    colorlinks=true,
    linkcolor=black,
    urlcolor=blue,
    citecolor=black,
    pdftitle={Relazione progetto Aurea -- OOP 2025/26},
    pdfauthor={Dapporto, Piraino, Conventi, Locatelli}
}
\usepackage{xcolor}
\usepackage{listings}
\lstset{
    basicstyle=\ttfamily\footnotesize,
    breaklines=true,
    columns=fullflexible,
    frame=single,
    showstringspaces=false,
    keywordstyle=\color{blue!70!black}\bfseries,
    commentstyle=\color{green!50!black}\itshape,
    stringstyle=\color{red!70!black},
    language=Java,
    captionpos=b,
    tabsize=2
}

% --- Indice e capitoli -------------------------------------------------------
\usepackage{titlesec}
\usepackage{parskip}
\setlength{\parskip}{0.4em}
\usepackage[titles]{tocloft}

% --- Frontespizio ------------------------------------------------------------
\title{%
    \vspace{-1.5cm}
    \rule{\linewidth}{0.5mm} \\[0.4cm]
    {\Huge \textbf{Aurea Mediocritas}} \\[0.2cm]
    {\Large Relazione del progetto di Programmazione ad Oggetti} \\ 
    \rule{\linewidth}{0.5mm} \\[0.4cm]
    {\large A.A. 2025/26 -- Universit\`a di Bologna}
}
\author{%
    Pietro Dapporto \\ \texttt{pietro.dapporto@studio.unibo.it}
    \and
    Nikita Piraino \\ \texttt{nikita.piraino@studio.unibo.it}
    \and
    Mario Conventi \\ \texttt{mario.conventi@studio.unibo.it}
    \and
    Federico Locatelli \\ \texttt{federico.locatelli3@studio.unibo.it}
}
\date{\today}

% =============================================================================
\begin{document}
% =============================================================================

\maketitle
\tableofcontents

% =============================================================================
\chapter{Analisi}
% =============================================================================
In questo capitolo si effettua l'analisi dei requisiti e del dominio applicativo
del progetto \emph{Aurea Mediocritas}, senza alcun riferimento al design o alle
tecnologie implementative.

% -----------------------------------------------------------------------------
\section{Descrizione e requisiti}
% -----------------------------------------------------------------------------
\emph{Aurea Mediocritas} \`e un'applicazione che riproduce le vicende di un
Rettore universitario impegnato a gestire un Ateneo bilanciando iniziative da
intraprendere e situazioni critiche da risolvere. Il gioco si ispira al gioco
di strategia a carte \emph{Reigns}: il giocatore impersona il Magnifico per la
durata di tre anni accademici (sei semestri) e deve arrivare al termine del
mandato mantenendo l'equilibrio fra i parametri che descrivono lo stato
dell'Universit\`a: \textbf{Finanze}, \textbf{Studenti}, \textbf{Docenti}
e \textbf{Reputazione}.

Prima di iniziare, il giocatore sceglie il proprio nome, la facolt\`a di
appartenenza e il livello di difficolt\`a della partita. A ogni turno riceve
una carta che descrive un evento (un'iniziativa proposta, una protesta
studentesca, un'opportunit\`a di finanziamento, eccetera) e deve scegliere
fra due alternative. Ciascuna scelta modifica i valori dei parametri di gioco.
La partita si conclude positivamente se il Rettore arriva al termine del
proprio mandato senza che alcun parametro raggiunga un valore critico;
si conclude negativamente nel momento in cui un qualsiasi parametro
raggiunge il livello minimo o massimo.

\paragraph{Requisiti funzionali}
\begin{itemize}
    \item Generazione di carte (eventi), ciascuna con un effetto associato 
          alla scelta del giocatore.
    \item Costruzione di un mazzo che modella la successione di eventi 
          della partita.
    \item Risposta del giocatore agli eventi tramite una scelta binaria 
          (approvazione o rifiuto).
    \item Aggiornamento dei parametri di gioco a seguito della scelta 
          effettuata, con intensità variabile in base alla difficoltà 
          selezionata.
    \item Gestione delle condizioni di fine partita: vittoria al termine 
          del mandato; sconfitta al raggiungimento dei valori critici di 
          un parametro (0 o 100).
    \item Sistema di difficoltà selezionabile (Facile, Normale, Difficile) 
          che influenza sia il peso dei danni applicati ai parametri sia 
          l'algoritmo di selezione delle carte.
    \item Sistema di carte figlie (\emph{follow-up}): alcune decisioni 
          programmano la comparsa di una carta correlata dopo un numero 
          definito di turni.
    \item Riepilogo inter-semestrale: al termine di ogni semestre viene 
          mostrata una schermata con i valori correnti dei quattro parametri.
    \item Popup del consigliere: quando un parametro attraversa una soglia 
          critica (inferiore al 10\% o superiore al 90\%), appare un avviso 
          contestuale con un messaggio narrativo.
    \item Schermata di login: il giocatore inserisce il proprio nome, 
          la facoltà e la difficoltà prima di iniziare la partita.
    \item Possibilità di riavviare la partita al termine senza chiudere 
          l'applicazione.
    \item Sistema audio: musica di sottofondo continua durante login 
          e partita, effetti sonori contestuali (scelta sulla carta, 
          comparsa del consigliere) e suoni dedicati per la vittoria 
          e la sconfitta.
\end{itemize}

\paragraph{Requisiti non funzionali}
\begin{itemize}
    \item L'interfaccia grafica deve essere reattiva e gradevole, con 
          animazioni fluide nelle transizioni di stato (ingresso carte, 
          aggiornamento parametri, eccetera).
\end{itemize}

% -----------------------------------------------------------------------------
\section{Modello del Dominio}
% -----------------------------------------------------------------------------
Lo svolgimento della partita è scandito da una sequenza di \emph{Carte}
prelevate da un \emph{Mazzo}. Ogni carta descrive un evento, è associata
a un \emph{Personaggio} che la presenta, e propone al giocatore due
\emph{Decisioni} alternative (approvazione o rifiuto). A ciascuna
decisione è associato un insieme di \emph{Effetti}, ognuno dei quali
specifica quanto e in quale direzione un parametro viene modificato. Ogni parametro
ha un nome, un livello numerico compreso fra 0 e 100, e uno stato
(attivo o esaurito). I quattro parametri sono: Finanze, Studenti,
Docenti e Reputazione.

Alcune decisioni attivano un meccanismo di \emph{Conseguenza}: Le scelte
prese, hanno conseguenza sulle situazioni che il rettore dovrà gestire 
durante la partita.

Il \emph{Tempo di gioco} è organizzato in turni raggruppati in
semestri; un mandato completo è composto da sei semestri di sei
turni ciascuno.

Prima di iniziare, il giocatore fornisce il livello di difficoltà a cui vuole giocare.


\begin{figure}[H]
    \centering
    \fbox{\includegraphics[width=0.8\textwidth]{UML-Piraino-Modello.png}}
    \caption{Modello del dominio di Aurea Mediocritas: entit\`a, attributi principali e relazioni.}
    \label{fig:domain-model}
\end{figure}

\paragraph{Aspetti impegnativi}
L'aspetto più complesso da modellare è l'algoritmo di selezione delle carte:
la successione non deve essere puramente casuale, ma deve evolvere in funzione
dello stato corrente dei parametri, della difficoltà selezionata e delle
conseguenze programmate dalle decisioni precedenti (follow-up).

% =============================================================================
\chapter{Design}
% =============================================================================
In questo capitolo si descrivono le strategie messe in campo per soddisfare i
requisiti identificati in fase di analisi: prima la visione architetturale
complessiva, poi gli aspetti di design di dettaglio pi\`u rilevanti svolti
singolarmente da ciascun membro del gruppo.

% -----------------------------------------------------------------------------
\section{Architettura}
% -----------------------------------------------------------------------------
L'architettura dell'applicazione segue il pattern
\textbf{Model-View-Controller (MVC)}, con una netta separazione delle
responsabilità fra i tre livelli. Il confine fra le componenti è definito
esclusivamente da interfacce: nessuno strato conosce l'implementazione
concreta degli altri, solo il contratto esposto.

\begin{itemize}
    \item Il \textbf{Model}, contenuto nel package \texttt{it.unibo.aurea.model}, 
        rappresenta lo stato della partita e la logica di gioco. Il suo principale 
        punto di ingresso è l'interfaccia \texttt{GameEngine}, che espone le operazioni 
        necessarie per avanzare nella partita, ottenere la carta corrente, applicare gli 
        effetti delle decisioni e consultare lo stato dei parametri dell'Ateneo. Attorno 
        a \texttt{GameEngine} collaborano componenti come \texttt{GameConfig}, per la 
        configurazione della partita, \texttt{GameClock}, per la progressione temporale, 
        \texttt{Deck}, per la gestione delle carte, e \texttt{Parameter}, per rappresentare 
        i quattro indicatori principali del gioco. Il Model non dipende dalla View concreta 
        né dal Controller: espone il proprio stato tramite interfacce e, per gli aggiornamenti 
        dei parametri, utilizza un meccanismo di notifica basato su observer.

    \item Il \textbf{Controller}, contenuto nel package \texttt{it.unibo.aurea.controller}, 
        ha come punto di ingresso l'interfaccia \texttt{GameController}. Esso coordina il flusso 
        dell'applicazione: avvia la partita tramite \texttt{startGame()}, riceve dalla View le 
        decisioni del giocatore tramite \texttt{makeDecision(boolean)}, le inoltra al Model 
        attraverso \texttt{GameEngine}, e recupera le informazioni necessarie per aggiornare 
        la rappresentazione grafica, ad esempio tramite \texttt{getCurrentParametersLevels()} 
        e \texttt{previewDecisionDeltas(boolean)}. In questo modo il Controller contiene la 
        logica di coordinamento, mentre non conserva direttamente lo stato del dominio, che 
        rimane responsabilità del Model.

    \item La \textbf{View}, contenuta nel package \texttt{it.unibo.aurea.view}, ha come punto di 
        ingresso l'interfaccia \texttt{GameView}. Essa si occupa esclusivamente della presentazione 
        dello stato del gioco e della raccolta degli input dell'utente. La View non accede direttamente 
        al Model, ma comunica con il sistema attraverso il \texttt{GameController}. Le operazioni come 
        \texttt{displayCard()}, \texttt{updateTime()} e l'aggiornamento dei parametri vengono invocate 
        dal Controller o attivate tramite il meccanismo di observer previsto per i cambiamenti dello 
        stato di gioco.
\end{itemize}

% -- inserire qui il diagramma UML architetturale --
\begin{figure}[H]
    \centering
    \fbox{\includegraphics[width=0.9\textwidth]{architettura_mvc.png}}
    \caption{Architettura complessiva del sistema: interfacce principali di Model, View e Controller.}
    \label{fig:architecture}
\end{figure}

Questa organizzazione consente di sostituire l'implementazione grafica senza modificare il Model. Ad 
esempio, una diversa interfaccia utente potrebbe implementare \texttt{GameView} e continuare a comunicare 
con il resto dell'applicazione attraverso \texttt{GameController}. Allo stesso modo, il Model rimarrebbe 
riutilizzabile anche in un'applicazione diversa che condivida la stessa logica di gioco, poiché non contiene 
riferimenti a elementi grafici o a classi specifiche della View.

% -----------------------------------------------------------------------------
\section{Design dettagliato}
% -----------------------------------------------------------------------------
Questa sezione approfondisce le principali scelte progettuali adottate da
ciascun membro del gruppo. Per ogni aspetto trattato vengono presentati:
il problema da risolvere nel contesto specifico del progetto, la soluzione
proposta con le motivazioni che l'hanno guidata, le alternative considerate
e scartate, uno schema UML di supporto e l'identificazione del pattern di
progettazione applicato, ove presente.

% --- 2.2.1 ---------------------------------------------------------
\subsection{Pietro Dapporto -- Modello delle carte e contenuti}
La mia responsabilità nel progetto riguarda la creazione e gestione delle carte.
Di seguito si riassumono i principali problemi incontrati con un accenno agli elementi creati
per risolverli.

\begin{table}[h]
\centering
\begin{tabular}{|l|l|l|}
\hline
\textbf{compito} & \textbf{nome interfaccia} & \textbf{nome implementazione} \\ \hline
Modello delle carte e contenuti & Card, Effect & CardImpl, EffectImpl, Decision \\ \hline
Scalare la creazione delle carte & - & Deck, CardsFile \\ \hline
Gestione carte figlie & FollowUp & FollowUpImpl \\ \hline
\end{tabular}
\end{table}

\subsubsection{Problema: modellazione delle carte di gioco}

Nel gioco ogni carta rappresenta una proposta sottoposta al giocatore da un determinato personaggio. Si deve modellare una carta che non sia un semplice insieme di dati, ma che permetta di rappresentare in modo chiaro le due possibili decisioni del giocatore e le relative conseguenze sui parametri della partita. 

\textbf{Soluzione} \\
La soluzione adottata consiste nel modellare la carta come un aggregato di oggetti più semplici, ciascuno con una responsabilità specifica. 
Una \texttt{Card} contiene un identificativo, una descrizione, il \texttt{Character} che la presenta e due oggetti \texttt{Decision}, rispettivamente associati al rifiuto e all'approvazione della proposta. 
Ogni \texttt{Decision} contiene il testo della risposta mostrata al giocatore e una collezione di \texttt{Effect}, che rappresentano le variazioni prodotte sui parametri del gioco.

In questo modo la carta non deve conoscere direttamente la logica di aggiornamento dei parametri, ma si limita a esporre le informazioni necessarie alla gestione della scelta. 

Infine, il flag \texttt{usage} permette di distinguere le carte già utilizzate da quelle ancora disponibili, evitando che la stessa carta venga riproposta più volte durante la partita.

\textbf{UML} \\
\begin{figure}[H]
    \centering
    \fbox{\includegraphics[width=0.70\textwidth]{problema_modellazione_carte.png}}
    \caption{Diagramma UML della carta e i suoi elementi}
    \label{fig:problema_modello_carte}
\end{figure}

\textbf{Pattern utilizzati} \\
In questa parte non è stato applicato un design pattern comportamentale in senso stretto. 
L'uso di interfacce come \texttt{Card} ed \texttt{Effect} consente comunque di mantenere il modello estendibile, poiché eventuali implementazioni alternative potrebbero essere introdotte senza modificare il codice che dipende dalle astrazioni.

\subsubsection{Problema: scalare la creazione delle carte}

La creazione manuale di un numero elevato di carte direttamente nel codice avrebbe reso il sistema poco leggibile, difficile da modificare e soggetto a ripetizioni. 
Il problema principale era quindi permettere l'aggiunta e la modifica delle carte in modo semplice e mantenibile, senza dover intervenire ogni volta sulla logica Java del model.

\textbf{Soluzione} \\
In una prima fase sono state considerate diverse soluzioni basate sulla costruzione diretta degli oggetti nel codice, ad esempio tramite il costruttore di \texttt{CardImpl}, tramite factory method o attraverso il pattern Builder. 
Queste soluzioni avrebbero però mantenuto la definizione delle carte troppo vicina al codice Java, rendendo meno immediata la lettura del contenuto di gioco e più costosa l'aggiunta di nuove carte.

Per questo motivo si è scelto di separare i dati dalla logica applicativa, spostando le informazioni delle carte in un file esterno in formato YAML. 
Il file contiene la descrizione dichiarativa delle carte, delle decisioni disponibili e degli effetti prodotti sui parametri del gioco. 
In questo modo il model non dipende dal modo in cui le carte sono scritte nel file, ma riceve oggetti già costruiti e coerenti con le primi astrazioni.

Il caricamento delle carte avviene all'interno della classe \texttt{Deck}, tramite metodi dedicati. 
Il file YAML viene deserializzato in una struttura intermedia rappresentata da \texttt{CardsFile}, che contiene una collezione di \texttt{CardDTO}: ciascuna contiene le decisioni \texttt{DecisionDTO} che a loro volta contengono i relativi \texttt{EffectDTO}. 
Durante questa fase, i dati letti dal file vengono convertiti dagli oggetti DTO agli oggetti effettivi del dominio, come \texttt{Card}, \texttt{Decision} ed \texttt{Effect}. 

Questa scelta migliora la mantenibilità perché nuove carte possono essere aggiunte o modificate intervenendo principalmente sul file di configurazione. 
La soluzione introduce però anche un livello di astrazione aggiuntivo: potenzialmente sarebbe possibile convertire le informazioni del file YAML in modelli di carta differenti, ma nel progetto attuale tutte le carte vengono ricondotte allo stesso modello interno. 

\textbf{UML} \\
\begin{figure}[H]
    \centering
    \fbox{\includegraphics[width=0.8\textwidth]{problema_scalabilità_carte.png}}
    \caption{Diagramma UML della creazione delle carte}
    \label{fig:problema_creazione_carte}
\end{figure}

\textbf{Pattern utilizzati} \\
In questa parte è stato utilizzato il pattern Builder per la creazione di \texttt{CardImpl} nel metodo \texttt{toCard()} dentro la classe \texttt{Deck}.
Il suo scopo è quello di convertire tutte le istanze di \texttt{CardDTO} in \texttt{CardImpl}.
Inoltre, per gestire il passaggio tra il file YAML e gli oggetti del model è stato utilizzato il concetto di DTO (\emph{Data Transfer Object}). 
Un DTO è un oggetto usato per trasferire dati tra due parti del sistema senza contenere logica di dominio significativa. 
In questo ambito agiscono quindi come classi ponte tra la rappresentazione esterna delle carte, definita nel file YAML, e la rappresentazione interna usata dal model. 

\subsubsection{Problema: gestione delle carte figlie}

Il gioco deve permettere che alcune decisioni producano conseguenze non immediate, facendo apparire nei turni successivi una carta collegata alla scelta effettuata. 
Inoltre la carta figlia deve apparire solo dopo una decisione su un'altra carta e non casualmente come una normale carta del mazzo principale.

\textbf{Soluzione} \\
Per implementare questo comportamento è stato introdotto il concetto di \texttt{FollowUp}, anch'esso definito tramite file YAML. 
Un \texttt{FollowUp} rappresenta una relazione logica tra una carta padre e una carta figlia, identificandole tramite i rispettivi \texttt{id} univoci. 
Oltre agli identificativi, il \texttt{FollowUp} specifica il trigger che attiva la relazione, cioè la scelta necessaria tra approvazione e rifiuto, e il numero di turni di ritardo dopo i quali la carta figlia potrà essere proposta.

Le carte figlie sono mantenute in un file YAML separato rispetto al mazzo principale. 
Durante il caricamento, il mazzo principale e le carte figlie vengono quindi convertiti in strutture dati distinte: il primo rappresenta le carte pescabili normalmente, mentre la seconda struttura contiene le carte disponibili solo come conseguenza di un \texttt{FollowUp}.

Quando il giocatore prende una decisione su una carta, il sistema verifica se esiste un \texttt{FollowUp} compatibile con l'id della carta e con la scelta effettuata. 
In caso positivo, la carta figlia viene programmata per comparire dopo il numero di turni indicato. 

\textbf{UML} \\
\begin{figure}[H]
    \centering
    \fbox{\includegraphics[width=0.8\textwidth]{problema_carte_figlie.png}}
    \caption{Diagramma UML del concetto di FollowUp}
    \label{fig:problema_arte_figlie}
\end{figure}

\textbf{Pattern utilizzati} \\
In questa parte non è stato applicato un design pattern classico in senso stretto. 
La soluzione si basa principalmente sulla separazione delle responsabilità e sulla rappresentazione esplicita delle relazioni tra carte tramite configurazione esterna.

% =============================================================================
% 2.2.2 Mario Conventi -- Motore di gioco, algoritmi adattivi e gestione conseguenze
% =============================================================================
\subsection{Mario Conventi -- Motore di gioco, algoritmi adattivi e gestione conseguenze}

La mia responsabilità principale all'interno del progetto riguarda la progettazione e lo sviluppo del nucleo computazionale del \emph{Model}: l'orchestratore delle regole di business \texttt{GameEngineImpl} e le sue componenti specializzate ad alta coesione. Il design si focalizza su tre pilastri algoritmici volti a trasformare un sistema a stati deterministico in un motore di gioco adattivo e configurabile.

\begin{table}[H]
\centering
\begin{tabular}{|l|l|l|}
\hline
\textbf{Compito} & \textbf{Interfaccia} & \textbf{Implementazione} \\ \hline
Motore di gioco (orchestratore) & \texttt{GameEngine} & \texttt{GameEngineImpl} \\ \hline
Selezione pesata adattiva delle carte & --- & \texttt{CardSelector} \\ \hline
Impostazioni di bilanciamento per difficoltà & --- & \texttt{DifficultySettings} \\ \hline
Coda delle carte figlie con ritardo & --- & \texttt{FollowUpQueue} \\ \hline
Contratto evento asincrono & \texttt{FollowUp} & \texttt{FollowUpImpl} \\ \hline
Test del motore con mazzo deterministico & --- & \texttt{GameEngineTest} (+ \texttt{TestDeck}) \\ \hline
\end{tabular}
\caption{Componenti realizzati da Mario Conventi.}
\label{tab:conventi-components}
\end{table}

\bigskip\noindent\rule{\linewidth}{0.4pt}\bigskip

% --- PILASTRO I -------------------------------------------------------------
\subsubsection{Pilastro I -- Algoritmo di estrazione a priorità rigida}

\begin{figure}[H]
    \centering
    \includegraphics[width=0.92\linewidth]{img/pilastro1_estrazione.png}
    \caption{Diagramma di sequenza UML (Pilastro I): flusso di estrazione a priorità rigida. Evidenzia le tre fasi sequenziali non intercambiabili del metodo \texttt{extractNextCard()}.}
    \label{fig:pilastro1}
\end{figure}

\textbf{Problema:} La progressione del flusso narrativo e delle meccaniche di gioco richiede l'estrazione continua di eventi dal mazzo. Tuttavia, la successione delle carte non può essere piatta: le carte generate come conseguenze storiche di decisioni passate (\emph{carte figlie}), una volta esaurito il loro tempo di incubazione, devono imporsi sul flusso standard del gioco con precedenza assoluta. Occorre garantire un punto di controllo centralizzato e sequenziale che eviti condizioni di errore logico o il salto accidentale di eventi forzati.

\textbf{Soluzione:} Il problema è stato risolto incapsulando la gerarchia di precedenza all'interno del metodo privato \texttt{extractNextCard()} in \texttt{GameEngineImpl}. L'algoritmo codifica una politica a priorità rigida non negoziabile strutturata in tre fasi sequenziali:

\begin{lstlisting}
private Card extractNextCard() {
    // Fase 1: Verifica della presenza di eventi forzati scaduti
    final Optional<Card> forcedCard = this.followUpQueue.pollForcedCard(this.deck);
    if (forcedCard.isPresent()) {
        return forcedCard.get();
    }
    // Fase 2: Solo se non c'e' carta forzata: incremento del contatore turni
    this.followUpQueue.updateTurns();
    
    // Fase 3: Selezione standard adattiva dal mazzo
    return this.cardSelector.selectNextCard(this.deck, this.parameters, this.difficultySettings);
}
\end{lstlisting}

L'ordine di esecuzione delle tre fasi non è casuale: il decremento dei contatori (\texttt{updateTurns()}) avviene rigorosamente \emph{dopo} il controllo delle scadenze e \emph{prima} della selezione ordinaria. Questo previene un diffuso errore di sfasamento (\emph{off-by-one}): se una carta figlia ha un tempo di ritardo pari a zero, essa viene erogata al turno successivo e non nello stesso istante della decisione. 

Alcune alternative considerate sono:
\begin{itemize}
    \item \emph{Aggiornamento dei turni all'inizio del metodo}: avrebbe anticipato la maturazione dei ritardi, facendo apparire le carte con un turno di anticipo rispetto alla semantica dichiarata nella configurazione.
    \item \emph{Gestione condizionale piatta}: avrebbe frammentato lo stato logico, esponendo il sistema a inconsistenze in caso di modifiche future.
\end{itemize}

\textbf{Eventuali pattern:} nessuno in senso stretto oltre alla strutturazione a priorità deterministica.

\bigskip\noindent\rule{\linewidth}{0.4pt}\bigskip

% --- PILASTRO II ------------------------------------------------------------
\subsubsection{Pilastro II -- Algoritmo di selezione pesata adattiva}

\begin{figure}[H]
    \centering
    \includegraphics[width=0.95\linewidth]{img/pilastro2_strategy.png}
    \caption{Diagramma delle classi UML (Pilastro II): pattern Strategy applicato alla selezione pesata. \texttt{CardSelector} incapsula l'algoritmo adattivo, modulato dal \emph{Parameter Object} \texttt{DifficultySettings}.}
    \label{fig:pilastro2}
\end{figure}

\textbf{Problema:} Un sistema basato sulla pura estrazione casuale rischia di rendere l'esperienza utente frustrante o ingiusta. Se un giocatore si trova in una situazione di grave crisi su un parametro, il motore deve calibrare la distribuzione delle carte per offrire opportunità di riequilibrio. Al contempo, si deve evitare un determinismo assoluto che eliminerebbe la rigiocabilità.

\textbf{Soluzione:} È stato progettato un algoritmo di \textbf{selezione casuale pesata} (\emph{weighted random selection}) nella componente \texttt{CardSelector}, che opera in quattro fasi logiche.

\textbf{Fase 1 -- Identificazione del parametro critico:} si individua il parametro con distanza minima dall'estremo (0 o 100), calcolata come $\min(\text{livello},\, 100 - \text{livello})$.

\textbf{Fase 2 -- Filtro di sicurezza letale:} si escludono le carte già usate, quelle figlie e le carte letali in entrambe le opzioni (\texttt{isLethalInBothOptions}), prevenendo scenari di morte certa inevitabile.

\textbf{Fase 3 -- Assegnazione dei pesi con boost adattivo:} ogni carta parte da un peso base (\texttt{DEFAULT\_WEIGHT = 10.0}); se aiuta il parametro critico, riceve un incremento dinamico:
\begin{lstlisting}
weight *= 1.0 + (NEUTRAL_DISTANCE - minDistance) / difficultySettings.getWeightDivisor();
\end{lstlisting}
La modalità \emph{Hard} imposta \texttt{weightDivisor = 500.0}, sopprimendo il boost e avvicinando la selezione alla casualità pura.

\textbf{Fase 4 -- Campionamento proporzionale (roulette wheel):} si genera un numero casuale in $[0,\,\text{totalWeight})$ e si scorre la lista cumulando i pesi fino a superare la soglia.

Alcune alternative considerate sono:
\begin{itemize}
    \item \emph{Selezione puramente casuale}: nessun bilanciamento adattivo; l'esperienza di gioco risulterebbe ingiusta nelle situazioni di crisi.
    \item \emph{Selezione deterministica}: la carta di aiuto verrebbe sempre estratta; il gioco perderebbe variabilità e rigiocabilità.
\end{itemize}

Pattern identificato Strategy e DTO (comportamentale e strutturale): \texttt{CardSelector} è una strategia di selezione intercambiabile: \texttt{GameEngineImpl} la usa tramite composizione e potrebbe sostituirla senza modificare il proprio codice. \texttt{DifficultySettings} agisce da \emph{Parameter Object} che incapsula i parametri di configurazione, preservando l'estendibilità (OCP). Per evitare l'uso di costanti hardcoded, i parametri come \texttt{weightDivisor} vengono caricati dinamicamente da un file di configurazione esterno (\texttt{difficulty.yaml}) tramite oggetti di trasferimento dati (\texttt{DifficultySettingsDTO} e \texttt{DifficultySettingsFile}), separando nettamente il bilanciamento del gioco dalla logica algoritmica.

\bigskip\noindent\rule{\linewidth}{0.4pt}\bigskip

% --- PILASTRO III -----------------------------------------------------------
\subsubsection{Pilastro III -- Gestione delle conseguenze e delle carte figlie}

\begin{figure}[H]
    \centering
    \includegraphics[width=0.90\linewidth]{img/pilastro3_followup.png}
    \caption{Diagramma delle classi UML (Pilastro III): infrastruttura di gestione dei FollowUp. Mostra la composizione \texttt{FollowUpQueue}~$\diamond$~\texttt{ActiveFollowUp} (0..*) e la realizzazione \texttt{FollowUpImpl}~$\dashrightarrow$~\texttt{FollowUp}.}
    \label{fig:pilastro3}
\end{figure}

\textbf{Problema:} Le decisioni del Rettore possono innescare eventi asincroni complessi, programmando la comparsa di una carta figlia (\emph{Follow-Up}) distanziata nel tempo da un numero di turni specifico (\texttt{delayTurn}). Il sistema deve supportare la registrazione atomica delle conseguenze, la maturazione passiva del tempo a ogni turno e l'estrazione prioritaria (Pilastro~I), preservando l'immutabilità del mazzo originario.

\textbf{Soluzione:} La gestione è stata centralizzata in \texttt{FollowUpQueue}, che incapsula una struttura dinamica basata sulla classe interna statica privata \texttt{ActiveFollowUp}. L'infrastruttura coordina due operazioni fondamentali.

\textbf{Registrazione (\texttt{registerConsequences}):} al momento della scelta binaria, la coda filtra tramite stream i contratti associati alla carta e all'esito (\emph{APPROVAL} o \emph{REFUSAL}) e li accoda con il loro \texttt{delayTurn} come tempo di attesa iniziale:
\begin{lstlisting}
deck.getAllFollowUps().stream()
    .filter(fu -> fu.getParentId().equals(parentId))
    .filter(fu -> fu.getTrigger() == actualOutcome)
    .forEach(fu -> eventQueue.add(new ActiveFollowUp(fu, fu.getDelayTurn())));
\end{lstlisting}

\textbf{Polling sicuro (\texttt{pollForcedCard}):} scansiona la coda per individuare gli eventi scaduti (\texttt{remainingTurns <= 0}) e li rimuove. Per prevenire \texttt{ConcurrentModificationException}, la rimozione durante la scansione avviene tramite \textbf{Iterator esplicito}:
\begin{lstlisting}
final Iterator<ActiveFollowUp> iterator = eventQueue.iterator();
while (iterator.hasNext()) {
    final ActiveFollowUp activeEvent = iterator.next();
    if (activeEvent.getRemainingTurns() <= 0) {
        iterator.remove();
        // ricerca della carta figlia nel pool tramite stream...
    }
}
\end{lstlisting}

Alcune alternative considerate sono:
\begin{itemize}
    \item \emph{\texttt{Queue<>} di Java}: non consente l'iterazione con rimozione selettiva basata su condizione interna.
    \item \emph{Ciclo \texttt{for-each} con lista separata di rimozioni}: più verboso e soggetto a errori; l'Iterator esplicito è la soluzione idiomatica Java per questo caso.
\end{itemize}

Pattern identificato Iterator esplicito (comportamentale): L'uso di un \texttt{Iterator} strutturato, invece del costrutto \texttt{for-each}, garantisce la totale correttezza in presenza di rimozioni durante la scansione, azzerando il rischio di eccezioni a runtime. Combinato con il \textbf{Test Double / Fake} (\texttt{TestDeck}), consente test unitari atomici e riproducibili al 100\%.

\bigskip\noindent\rule{\linewidth}{0.4pt}\bigskip

% --- 2.2.3 ---------------------------------------------------------
\subsection{Federico Locatelli -- Parametri reattivi, architettura della View e ciclo di vita}

La mia responsabilità copre tre livelli del sistema: la gestione reattiva
dei parametri nel Model, l'intero layer View con le sue animazioni e
componenti interattivi, e l'estensione del Controller con nuovi metodi
di anteprima e accesso allo stato. Di seguito una panoramica dei
componenti realizzati.

\begin{table}[h]
\centering
\begin{tabular}{|l|l|l|}
\hline
\textbf{Compito} & \textbf{Interfaccia} & \textbf{Implementazione} \\ \hline
Parametro (sola lettura, View layer) & \texttt{ParameterView} & \texttt{ParameterImpl} \\ \hline
Parametro (mutabile, Model layer) & \texttt{Parameter} & \texttt{ParameterImpl} \\ \hline
Observer dei parametri & \texttt{ParameterObserver} & lambda / method reference \\ \hline
Tipo parametro con comportamento & --- & \texttt{ParameterType} (enum) \\ \hline
Identità giocatore & --- & \texttt{PlayerInfo} (record) \\ \hline
Contratto del Controller & \texttt{GameController} & aggiunti a \texttt{GameControllerImpl} \\ \hline
View principale & \texttt{GameView} & \texttt{GameViewJavaFXImpl} \\ \hline
Carta interattiva swipeable & --- & \texttt{CardPanel} \\ \hline
Icona parametro animata & --- & \texttt{ParameterIconView} \\ \hline
Schermata di login & --- & \texttt{LoginScene} \\ \hline
Popup del consigliere & --- & \texttt{CounsellorPopup} \\ \hline
Overlay di fine partita (in collaborazione) & --- & \texttt{EndgameOverlay} \\ \hline
Gestione audio di gioco & --- & \texttt{AudioManager} \\ \hline
Stile visivo dell'applicazione & --- & \texttt{styles.css} \\ \hline
\end{tabular}
\caption{Componenti realizzati da Federico Locatelli.}
\label{tab:locatelli-components}
\end{table}

\bigskip
\subsubsection{Pattern Observer: parametri reattivi}
\smallskip

\paragraph{Problema:}
La View deve aggiornarsi ogni volta che il valore di un parametro cambia,
senza che il Model conosca l'esistenza della View. Un accoppiamento diretto
fra le due componenti violerebbe il principio di separazione delle
responsabilità dell'architettura MVC e renderebbe impossibile sostituire
la View senza modificare il Model.

\paragraph{Soluzione:}
È stato adottato il pattern \textbf{Observer}. L'interfaccia
\texttt{ParameterObserver} è dichiarata \texttt{@FunctionalInterface}:
definisce un unico metodo \texttt{onParameterChanged(ParameterType, int)}
che viene invocato ogni volta che il livello del parametro cambia.
\texttt{ParameterImpl} funge da \emph{observable}: mantiene una lista di
osservatori registrati e li notifica al termine di ogni \texttt{modify()},
usando una copia difensiva della lista (\texttt{List.copyOf}) per prevenire
eccezioni di modifica concorrente.

La registrazione dell'osservatore avviene nel costruttore di
\texttt{GameControllerImpl}, dove la View viene collegata al Model tramite
una \emph{method reference}:

\begin{lstlisting}
parametersMap.values()
    .forEach(p -> p.addObserver(this.view::updateSingleParameter));
\end{lstlisting}

In questo modo il Model non conosce la View: conosce solo l'interfaccia
\texttt{ParameterObserver}, che può essere implementata da qualsiasi
oggetto o, come in questo caso, da una lambda o method reference.

\paragraph{Alternative considerate:}
\begin{itemize}
    \item \emph{Polling}: la View interroga periodicamente il Model per
        verificare se i valori sono cambiati. Scartato perché introduce
        latenza, spreco di cicli CPU e accoppiamento temporale.
    \item \emph{Chiamata diretta View $\to$ Model}: il Controller aggiorna
        la View dopo ogni decisione. Parzialmente adottato per
        l'inizializzazione, ma insufficiente per aggiornamenti
        intermedi (es. effetti concatenati su più parametri).
\end{itemize}

\begin{figure}[H]
    \centering
    \fbox{\includegraphics[width=0.6\textwidth]{UML-Locatelli-Observer.png}}
    \caption{Pattern Observer applicato ai parametri di gioco.}
    \label{fig:observer}
\end{figure}

\paragraph{Pattern identificato -- Observer:}
\texttt{ParameterImpl} è l'\emph{observable}; \texttt{ParameterObserver}
è l'interfaccia dell'\emph{observer}; \texttt{GameViewJavaFXImpl}
(\texttt{updateSingleParameter}) è l'\emph{observer concreto}, registrato
tramite method reference dal Controller.

\bigskip\noindent\rule{\linewidth}{0.4pt}\bigskip

\bigskip
\subsubsection{Interface Segregation: separazione tra lettura e scrittura dei parametri}
\smallskip

\paragraph{Problema:}
Il Model espone i parametri sia alla View (che deve solo leggerli per
aggiornare l'interfaccia) sia al Controller (che deve modificarli
applicando gli effetti delle carte). Esporre un'unica interfaccia con
tutti i metodi, incluso \texttt{modify()}, avrebbe consentito alla View
di alterare direttamente lo stato del gioco, violando i confini
dell'architettura MVC.

\paragraph{Soluzione:}
È stato applicato l'\textbf{Interface Segregation Principle}: i metodi
del parametro sono stati suddivisi in due interfacce distinte con
responsabilità diverse.

\texttt{ParameterView} espone esclusivamente i metodi di consultazione:
\texttt{getLevel()}, \texttt{isAlive()}, \texttt{getName()},
\texttt{getDeathReason()} e \texttt{addObserver()}. Questo è il tipo
che la View e il Controller utilizzano per leggere lo stato dei parametri
senza rischiare modifiche accidentali.

\texttt{Parameter} estende \texttt{ParameterView} aggiungendo il solo
metodo \texttt{modify(int delta)}, che altera il livello del parametro.
Questo tipo è utilizzato esclusivamente all'interno del Model.

La distinzione è resa esplicita nel Controller, dove la mappa dei
parametri è tipizzata come \texttt{Map<ParameterType, ParameterView>}:
il Controller legge i livelli e registra gli observer, ma non chiama
mai \texttt{modify()} direttamente --- delega questa responsabilità
al Model tramite \texttt{model.applyEffects()}.

\paragraph{Alternative considerate:}
\begin{itemize}
    \item \emph{Interfaccia unica}: esporre tutti i metodi in una sola
        interfaccia. Più semplice, ma permette alla View di chiamare
        \texttt{modify()} per errore, introducendo bug difficili da
        diagnosticare.
    \item \emph{Copia difensiva al getter}: restituire una copia dei
        parametri invece di riferimenti diretti. Scartato perché rompe
        il collegamento con il pattern Observer --- gli observer registrati
        sulla copia non riceverebbero le notifiche del parametro originale.
\end{itemize}

% -- UML ISP --
\begin{figure}[H]
    \centering
    \fbox{\includegraphics[width=0.3\textwidth]{UML-Locatelli-Interface-Segregation.png}}
    \caption{Interface Segregation applicata ai parametri di gioco.}
    \label{fig:isp}
\end{figure}

\paragraph{Principio identificato -- Interface Segregation (ISP):}
\texttt{ParameterView} è l'interfaccia \emph{ristretta} usata da View
e Controller; \texttt{Parameter} è l'interfaccia \emph{completa} usata
dal Model; \texttt{ParameterImpl} implementa entrambe.

\bigskip\noindent\rule{\linewidth}{0.4pt}\bigskip

\bigskip
\subsubsection{Type-Safe Enum con comportamento: ParameterType}
\smallskip

\paragraph{Problema:}
I quattro parametri del gioco devono essere identificati in modo
univoco e presentati all'utente con un nome leggibile. Una soluzione
ingenua avrebbe usato stringhe letterali (\texttt{"finances"},
\texttt{"students"}, ecc.) o avrebbe applicato trasformazioni come
\texttt{name().toLowerCase()} direttamente nel codice della View,
introducendo dipendenze fragili dal nome interno dell'enum e
duplicazione della logica di presentazione.

\paragraph{Soluzione:}
È stato adottato il pattern \textbf{Type-Safe Enum con comportamento}:
l'enum \texttt{ParameterType} non è un semplice elenco di etichette,
ma un oggetto che conosce come presentarsi. Ogni costante porta con
sé un campo \texttt{displayName} inizializzato nel costruttore e
accessibile tramite il metodo \texttt{getDisplayName()}:

\begin{lstlisting}
public enum ParameterType {
    FINANCES("Finances"),
    STUDENTS("Students"),
    PROFESSORS("Professors"),
    REPUTATION("Reputation");

    private final String displayName;

    ParameterType(final String displayName) {
        this.displayName = displayName;
    }

    public String getDisplayName() {
        return displayName;
    }
}
\end{lstlisting}

In questo modo la View chiama \texttt{type.getDisplayName()} invece
di manipolare \texttt{type.name()}: il codice di presentazione è
disaccoppiato dall'identificatore interno dell'enum. Se in futuro
si volesse rinominare la costante \texttt{FINANCES} in \texttt{MONEY},
il nome mostrato all'utente rimarrebbe invariato senza toccare la View.

\paragraph{Alternative considerate:}
\begin{itemize}
    \item \emph{Metodo statico di mappatura}: una \texttt{Map} o uno
        \texttt{switch} separato che traduce l'enum in stringa. Scartato
        perché distribuisce la logica di presentazione in un punto
        lontano dalla definizione del tipo.
    \item \emph{Override di toString()}: restituire il display name
        da \texttt{toString()}. Scartato perché \texttt{toString()} è
        un metodo generico con semantica ambigua; un metodo dedicato
        \texttt{getDisplayName()} rende l'intenzione esplicita.
\end{itemize}

% -- UML Type-Safe Enum --
\begin{figure}[H]
    \centering
    \fbox{\includegraphics[width=0.8\textwidth]{UML-Locatelli-Enum.png}}
    \caption{Type-Safe Enum con comportamento applicato a ParameterType.}
    \label{fig:enum}
\end{figure}

\paragraph{Pattern identificato -- Type-Safe Enum con comportamento:}
\texttt{ParameterType} non è solo un'etichetta ma un oggetto con
logica di presentazione incapsulata. Ogni costante è \emph{self-describing}:
sa come mostrarsi alla View senza richiedere logica esterna.

\bigskip\noindent\rule{\linewidth}{0.4pt}\bigskip

\bigskip
\subsubsection{Pattern Strategy: comportamenti iniettati in CardPanel}
\smallskip

\paragraph{Problema:}
\texttt{CardPanel} gestisce l'interazione fisica con la carta
(trascinamento, rotazione, volo), ma deve anche: notificare il
Controller quando il giocatore prende una decisione, richiedere
un'anteprima dei parametri colpiti durante il trascinamento, e
segnalare quando l'anteprima termina. Implementare questi comportamenti
direttamente in \texttt{CardPanel} avrebbe creato una dipendenza diretta
verso il Controller, accoppiando il componente visuale alla logica
di gioco e rendendo impossibile testarlo in isolamento.

\paragraph{Soluzione:}
È stato adottato il pattern \textbf{Strategy} tramite interfacce
funzionali Java: i tre comportamenti variabili sono rappresentati
da campi di tipo \texttt{BiConsumer}, \texttt{Function} e
\texttt{Runnable}, iniettati dall'esterno tramite metodi setter:

\begin{lstlisting}
private BiConsumer<Card, Boolean> onDecision;
private Function<Boolean, Set<ParameterType>> previewProvider;
private Runnable onPreviewEnd;
\end{lstlisting}

\texttt{CardPanel} non conosce il Controller: si limita a invocare
le strategie ricevute nei momenti opportuni. La configurazione
avviene in \texttt{GameViewJavaFXImpl}, che inietta le implementazioni
concrete tramite method reference verso il Controller:

\begin{lstlisting}
cardPanel.setOnDecision((card, approved) ->
    controller.makeDecision(approved));
cardPanel.setPreviewProvider(
    controller::previewDecisionDeltas);
cardPanel.setOnPreviewEnd(this::resetHighlights);
\end{lstlisting}

Questo approccio permette di sostituire i comportamenti senza
modificare \texttt{CardPanel} e di testarlo con implementazioni
fittizie (stub) senza coinvolgere il Controller reale.

\paragraph{Alternative considerate:}
\begin{itemize}
    \item \emph{Dipendenza diretta da GameController}: passare il
        Controller come parametro al costruttore di \texttt{CardPanel}.
        Scartato perché lega il componente visuale a una specifica
        interfaccia del Controller, riducendo la riusabilità.
    \item \emph{Listener ad hoc con interfacce dedicate}: definire
        interfacce \texttt{DecisionListener}, \texttt{PreviewListener}
        ecc. Più verboso senza vantaggi concreti rispetto all'uso
        diretto delle interfacce funzionali di Java.
\end{itemize}

% -- UML Strategy --
\begin{figure}[H]
    \centering
    \fbox{\includegraphics[width=0.8\textwidth]{UML-Locatelli-Strategy.png}}
    \caption{Pattern Strategy applicato ai comportamenti di CardPanel.}
    \label{fig:strategy}
\end{figure}

\paragraph{Pattern identificato -- Strategy:}
\texttt{CardPanel} è il \emph{context}; \texttt{BiConsumer},
\texttt{Function} e \texttt{Runnable} sono le \emph{strategie};
le lambda e method reference iniettate da \texttt{GameViewJavaFXImpl}
sono le \emph{strategie concrete}.

\bigskip\noindent\rule{\linewidth}{0.4pt}\bigskip

\bigskip
\subsubsection{Inversion of Control: gestione del ciclo di vita}
\smallskip

\paragraph{Problema:}
Al termine di una partita il giocatore può scegliere di ricominciare
da capo. La View, che gestisce la schermata di fine partita, conosce
il momento in cui il giocatore preme ``Reign Again'', ma non dovrebbe
sapere \emph{come} riavviare il gioco: farlo significherebbe che la
View conosce la struttura dell'intera applicazione (Model, Controller,
login), violando la separazione delle responsabilità MVC.

\paragraph{Soluzione:}
È stato applicato il principio di \textbf{Inversion of Control} tramite
callback: la View riceve un oggetto \texttt{Runnable onRestart} nel
costruttore, senza sapere cosa fa. Quando il giocatore preme il tasto
di riavvio, la View chiude gli stage aperti e invoca il callback,
delegando completamente la responsabilità a chi lo ha fornito.

Il proprietario del ciclo di vita è \texttt{Main}, che definisce il
callback come una method reference verso il proprio metodo
\texttt{openLogin()}:

\begin{lstlisting}
final Runnable onRestart = () ->
    Platform.runLater(this::openLogin);
final GameView view =
    new GameViewJavaFXImpl(onRestart);
\end{lstlisting}

Quando il callback viene eseguito, \texttt{Main} chiude il gioco
corrente e apre una nuova \texttt{LoginScene}, ricreando l'intero
stack MVC con i nuovi dati del giocatore. La View non conosce
nulla di questo flusso: esegue un \texttt{Runnable} e basta.

\paragraph{Alternative considerate:}
\begin{itemize}
    \item \emph{Riferimento diretto a Main}: passare l'istanza di
        \texttt{Main} alla View. Scartato perché crea una dipendenza
        ciclica e concettualmente errata: la View non dovrebbe
        conoscere il punto di ingresso dell'applicazione.
    \item \emph{Riavvio tramite Platform.exit() e riapertura}:
        chiudere la JVM e riaprire l'applicazione. Scartato perché
        non è trasparente all'utente e non è un riavvio autentico
        all'interno della stessa sessione.
\end{itemize}

% -- UML IoC --
\begin{figure}[H]
    \centering
    \fbox{\includegraphics[width=0.8\textwidth]{UML-Locatelli-InversionOfControl.png}}
    \caption{Inversion of Control per la gestione del riavvio.}
    \label{fig:ioc}
\end{figure}

\paragraph{Principio identificato -- Inversion of Control:}
\texttt{Main} è il \emph{proprietario del ciclo di vita};
\texttt{GameViewJavaFXImpl} è il \emph{consumatore del callback};
il \texttt{Runnable} è il \emph{contratto minimo} che disaccoppia
i due senza che la View conosca Main.

\bigskip\noindent\rule{\linewidth}{0.4pt}\bigskip

\bigskip
\subsubsection{Value Object: PlayerInfo come record immutabile}
\smallskip

\paragraph{Problema:}
I dati inseriti dal giocatore nella schermata di login --- nome del
Rettore, facoltà e difficoltà --- devono essere raccolti, validati
e trasportati attraverso il sistema fino al Controller, che li
utilizza per personalizzare l'esperienza di gioco. Una soluzione
basata su un oggetto mutabile avrebbe esposto il rischio di
modifiche accidentali durante il transito, e la validazione
sarebbe potuta essere dimenticata in qualunque punto del codice.

\paragraph{Soluzione:}
È stato adottato il pattern \textbf{Value Object} tramite il
costrutto \texttt{record} di Java~21. \texttt{PlayerInfo} è
dichiarato come record con tre componenti: \texttt{rectorName},
\texttt{faculty} e \texttt{difficulty}. Il costruttore compatto
applica la validazione \emph{fail-fast} su tutti i campi:

\begin{lstlisting}
public record PlayerInfo(
        String rectorName,
        String faculty,
        Difficulty difficulty) {

    public PlayerInfo {
        Objects.requireNonNull(rectorName,
            "rectorName must not be null");
        Objects.requireNonNull(faculty,
            "faculty must not be null");
        Objects.requireNonNull(difficulty,
            "difficulty must not be null");
    }
}
\end{lstlisting}

Il compilatore genera automaticamente \texttt{equals()},
\texttt{hashCode()} e \texttt{toString()} coerenti con i valori
dei campi, e i campi sono implicitamente \texttt{final}: una volta
creato, un \texttt{PlayerInfo} non può essere modificato. Questo
garantisce che i dati del giocatore siano consistenti in ogni punto
del sistema che li utilizza.

\paragraph{Alternative considerate:}
\begin{itemize}
    \item \emph{Classe mutabile con getter e setter}: più flessibile
        ma priva delle garanzie di immutabilità. La validazione
        dovrebbe essere replicata in ogni setter, aumentando
        il rischio di inconsistenza.
    \item \emph{Parametri separati}: passare nome, facoltà e
        difficoltà come argomenti separati al costruttore del
        Controller. Scartato perché aumenta l'arità dei metodi
        e non raggruppa semanticamente i dati correlati.
\end{itemize}

% -- UML Value Object --
\begin{figure}[H]
    \centering
    \fbox{\includegraphics[width=0.8\textwidth]{UML-Locatelli-ValueObject.png}}
    \caption{Value Object \texttt{PlayerInfo} come Java record.}
    \label{fig:playerinfo}
\end{figure}

\paragraph{Pattern identificato -- Value Object:}
\texttt{PlayerInfo} è immutabile, si confronta per valore e
trasporta dati validati. Il costrutto \texttt{record} di Java~21
fornisce queste garanzie con codice minimo e senza boilerplate.

\bigskip\noindent\rule{\linewidth}{0.4pt}\bigskip

\bigskip
\subsubsection{Ulteriori scelte di design}
\smallskip
Oltre ai pattern principali descritti nelle sezioni precedenti,
alcune scelte progettuali minori meritano una menzione per il loro
impatto sulla qualità e manutenibilità del codice.

\bigskip
\paragraph{Factory Method in GameViewJavaFXImpl:}
\smallskip
La costruzione dei componenti visivi della View è delegata a
metodi privati dedicati: \texttt{buildTopSection()},
\texttt{buildCenterSection()}, \texttt{buildBottomSection()},
\texttt{buildCloseButton()}. Ogni metodo è responsabile di
costruire e restituire un nodo JavaFX autocontenuto. Questo
approccio riduce la complessità del metodo principale
\texttt{buildAndShowStage()} e permette di modificare un
componente senza impattare gli altri.

\bigskip
\paragraph{Utility Class -- CounsellorPopup:}
\smallskip
\texttt{CounsellorPopup} è progettata come \emph{utility class}:
costruttore privato, nessun campo di istanza, unico metodo
pubblico statico \texttt{show(String message)}. La scelta è
motivata dalla natura \emph{stateless} del popup: ogni
invocazione costruisce un nuovo stage modale, lo mostra e
lo chiude senza mantenere stato tra una chiamata e l'altra.

\bigskip
\paragraph{Fail-Fast e Defensive Copy in ParameterImpl:}
\smallskip
Il costruttore di \texttt{ParameterImpl} e il metodo
\texttt{addObserver()} applicano \texttt{Objects.requireNonNull}
per rigettare immediatamente input nulli, evitando errori
latenti difficili da diagnosticare. Il metodo
\texttt{notifyObservers()} itera su una copia difensiva
della lista degli osservatori (\texttt{List.copyOf}), rendendo
il sistema robusto rispetto a modifiche concorrenti della lista
durante la notifica.

\bigskip
\paragraph{Gestione del threading in JavaFX:}
\smallskip
Tutti i metodi della View che modificano nodi della scena
grafica (\texttt{displayCard()}, \texttt{updateSingleParameter()},
\texttt{showVictory()}, ecc.) delegano l'esecuzione al thread
JavaFX tramite \texttt{Platform.runLater()}. Questo garantisce
che gli aggiornamenti dell'interfaccia avvengano sempre sul
thread corretto, prevenendo eccezioni di accesso concorrente
tipiche delle applicazioni JavaFX multi-thread.

\bigskip
\paragraph{Sistema audio: AudioManager e gestione del ciclo di vita dei player:}
\smallskip
\texttt{AudioManager} è progettata come \emph{utility class}
(costruttore privato, soli metodi statici) e centralizza tutta la
logica audio del gioco. Distingue due tipologie di suono: la musica
di sottofondo, gestita da un \texttt{MediaPlayer} in loop infinito
messo in pausa durante i suoni speciali e ripreso automaticamente
tramite \texttt{setOnEndOfMedia()}, e gli effetti brevi come lo swipe
della carta, affidati ad \texttt{AudioClip} per la bassa latenza.
Una scelta progettuale non banale riguarda la conservazione dei
riferimenti ai player di vittoria e sconfitta in un campo statico:
senza tale riferimento il \emph{garbage collector} di Java
interromperebbe la riproduzione prima del termine del suono, dato
che il \texttt{MediaPlayer} locale verrebbe deallocato all'uscita
dal metodo.

% --- 2.2.4 ---------------------------------------------------------
\subsection{Nikita Piraino -- Game loop, tempo e condizioni di fine partita}

La parte di progetto di mia responsabilità copre la logica di stato della partita e il ciclo temporale del gioco:

\begin{table}[h]
\centering
\begin{tabular}{|l|l|l|}
\hline
\textbf{compito} & \textbf{nome interfaccia} & \textbf{nome implementazione} \\ \hline
Configurazione partita & GameConfig & GameConfigImpl \\ \hline
Orchestratore logica di gioco & GameEngine & GameEngineImpl \\ \hline
Macchina a stati della partita & GameEngine & GameEngineImpl \\ \hline
Separazione accessi parametri & ParameterView & ParameterImpl \\ \hline
Riepilogo parametri partita & Report & ReportImpl \\ \hline
\end{tabular}
\end{table}

\subsubsection{Pattern Static Factory Method in GameConfigFactory}

\noindent\textbf{Problema} Sono previste più versioni di configurazione del gioco: una standard (6 carte per semestre, 6 semestri), una ridotta (2 carte, 2 semestri) pensata per facilitare i test intermedi e di fine partita. C'è anche il desiderio di implementare future versioni alternative di partite come ad esempio la modalità ``apocalisse''.
Occorre quindi controllare l'istanziamento della configurazione, vincolando gli utilizzatori a versioni standardizzate di partita, facilmente selezionabili tramite metodi pubblici e statici.

\noindent\noindent\textbf{Soluzione} Il pattern Static Factory Method è articolato in due classi distinte:

\begin{enumerate}
\item Costruttore package-private. Il costruttore \texttt{GameConfigImpl} è dichiarato con visibilità package-private: impedisce l'istanziazione diretta da parte di client esterni al package e concentra il controllo della validità dei parametri in un unico punto.
\item Classe factory dedicata. Le costanti di configurazione e i factory method sono spostati in \texttt{GameConfigFactory}, classe \texttt{final} con costruttore privato.

\item Factory method pubblici e statici. \texttt{createStandard(Difficulty)} produce la configurazione normale, \texttt{createShort(Difficulty)} la ridotta per testing; ogni ulteriore configurazione può essere aggiunta come nuovo metodo in \texttt{GameConfigFactory} senza modificare nessun altro codice se non la chiamata del metodo.
\end{enumerate}

Alcune alternative:
\begin{itemize}
\item Factory method nella stessa classe (\texttt{GameConfigImpl}): versione precedentemente adottata, nella quale i factory method e le costanti risiedevano direttamente nell'implementazione. Garantiva anch'essa information hiding, ma accoppiava responsabilità di creazione e di rappresentazione della configurazione nella stessa classe, risultando meno flessibile nell'eventualità di aggiungere configurazioni senza toccare l'implementazione.
\item Interfaccia come unico punto di accesso: sarebbe stato possibile garantire un information hiding totale evitando l'importazione di un'implementazione come \texttt{GameConfigImpl}; questo però avrebbe richiesto di rendere pubblico il costruttore dell'implementazione (Java non supporta la visibilità di costruttori tra i package /model e /model/api) e di spostare costanti e logica nell'interfaccia stessa, soluzione ritenuta poco elegante dal punto di vista architetturale e semantico.
\item Abstract Factory: avrebbe garantito information hiding, ma sarebbe risultato troppo verboso per varianti la cui differenza è solo di qualche variabile numerica o future aggiunte di difficoltà.
\end{itemize}

\begin{figure}[H]
    \centering
    \fbox{\includegraphics[width=0.8\textwidth]{UML-Piraino-Factory.png}}
    \caption{Pattern Static Factory Method in GameConfigFactory}
    \label{fig:factory}
\end{figure}

\textbf{Pattern identificato} Static Factory Method (creazionale): i metodi \texttt{createStandard()} e \texttt{createShort()} sono le factory; \texttt{GameConfig} è il tipo prodotto; \texttt{GameConfigFactory} è il punto di accesso pubblico.


\subsubsection{Orchestrazione del modello in GameEngineImpl}

\noindent\textbf{Problema} La logica di partita richiede di coordinare più oggetti collaboranti: configurazione, orologio, mazzo, parametri. Occorre un punto di ingresso unico per il controller, garantendo immutabilità strutturale del modello dopo la costruzione e fail-fast sulle dipendenze nulle.


\noindent\textbf{Soluzione} \texttt{GameEngineImpl} implementa \texttt{GameEngine} ed è la classe centrale del sottosistema model. Coordina tre dipendenze: la configurazione (\texttt{GameConfig}), l'orologio (\texttt{GameClock}) e il mazzo (\texttt{Deck}).

Il costruttore applica la Dependency Injection ricevendo \texttt{GameConfig}, \texttt{GameClock} e \texttt{Deck} come argomenti:

\begin{lstlisting}
public GameEngineImpl(final GameConfig config, final GameClock gameClock, final Deck deck) {
    this.config = Objects.requireNonNull(config, "GameConfig cannot be null");
    this.gameClock = Objects.requireNonNull(gameClock, "GameClock cannot be null");
    this.deck = Objects.requireNonNull(deck, "Deck cannot be null");
}
\end{lstlisting}

La chiamata a \texttt{Objects.requireNonNull} applica il principio fail-fast: se il client passa null, una NullPointerException viene sollevata immediatamente nel costruttore con un messaggio esplicativo, evitando errori latenti che si manifesterebbero solo durante l'esecuzione.

\begin{figure}[H]
    \centering
    \fbox{\includegraphics[width=0.8\textwidth]{UML-Piraino-Orchestrazione.png}}
    \caption{Orchestrazione del modello in GameEngineImpl}
    \label{fig:orchestrazione}
\end{figure}

un'alternativa:
\begin{itemize}
    \item utilizzare singleton pattern: in caso di gioco offline garantisce che ci sia un solo motore di un videogioco, però considerando il riuso del codice non sarebbe utilizzabile per costruire un server per partite online (dove, presumibilmente,
si gestiscono parallelamente più partite).
\end{itemize}

\subsubsection{State Machine in getGameState()}

\textbf{Problema:} Lo stato di una partita può essere in corso, vinta o persa; queste tre alternative sono mutuamente esclusive e dipendono dallo stato corrente dei parametri e dell'orologio. La consultazione dello stato è frequente (a ogni interazione del controller con il modello), per cui un campo di stato esplicito introdurrebbe rischi di disallineamento fra i dati interni e la valutazione restituita.

\noindent\textbf{Soluzione:} Il metodo \texttt{getGameState()} nella classe \texttt{GameEngineImpl} avvia una processo di identificazione dello stato, che poi restituirà uno dei valori dell'enumerazione GameState: RUNNING, WON, LOST.

\begin{lstlisting}
public GameState getGameState() {
    if (!areAllParametersAlive()) {
        return GameState.LOST;
    }
    if (gameClock.isTimeFinished()) {
        return GameState.WON;
    }
    return GameState.RUNNING;
}
\end{lstlisting}

L'ordine dei controlli è semanticamente obbligatorio: il controllo LOST (morte di un parametro) precede WON (tempo esaurito) perché nell'ultimo turno di gioco entrambe le condizioni possono essere simultaneamente vere; in tal caso la morte di un parametro deve prevalere sulla vittoria per tempo.

Lo stato viene calcolato ad ogni chiamata: l'approccio è stateless, elimina il rischio di stato inconsistente tra le transizioni e non richiede aggiornamenti espliciti dopo ogni operazione.

\begin{figure}[H]
    \centering
    \fbox{\includegraphics[width=0.8\textwidth]{UML-Piraino-GameState.png}}
    \caption{State Machine in GameEngineImpl}
    \label{fig:gamestate}
\end{figure}

\noindent Alcune alternative sono:
\begin{itemize}
    \item restituire una stringa al posto degli elementi dell'enum: non idoneo con gli standard di programmazione           moderni
    \item nelle condizioni non richiamare delle funzioni ma mettere il loro contenuto direttamente nel metodo \texttt{getGameState()}: supporta meno per necessità future i principi Don't Repeat Yourself e Single Responsibility.
\end{itemize}

\subsubsection{Separazione accessi Parameter e ParameterView}

\noindent\textbf{Problema} Esistono determinati oggetti che necessitano dell'accesso diretto agli oggetti parametri in lettura.

\noindent\textbf{Soluzione} La soluzione adottata separa la lettura dalla scrittura attraverso due interfacce distinte, rispettando pure l'Interface Segregation Principle:

\begin{enumerate}
\item Interfaccia di sola lettura. \texttt{ParameterView} dichiara esclusivamente i metodi di interrogazione del parametro, rappresentando una vista immutabile sicura da esporre all'esterno del Model.
\item Interfaccia mutabile. \texttt{Parameter} estende \texttt{ParameterView} aggiungendo il solo metodo di modifica.
\end{enumerate}

\begin{figure}[H]
    \centering
    \fbox{\includegraphics[width=0.8\textwidth]{UML-Piraino-ParameterView.png}}
    \caption{Separazione di Parameter e ParameterView}
    \label{fig:parameterview}
\end{figure}

Alcune alternative sono:
\begin{itemize}
\item Esporre direttamente \texttt{Parameter}: avere un metodo \texttt{getParameters()} che restituiva il tipo mutabile e invece un metodo che \texttt{getCopyOfParameters} che restituiva una copia difensiva e che doveva essere usata da chiunque non avesse bisogno di modificare il valore dei parametri. Risultava più semplice, ma lasciava a chiunque la possibilità di chiamare il metodo di modifica, rendendo l'incapsulamento dello stato puramente convenzionale e non garantito dal compilatore.
\item restituire una struttura dati contenente i valori dei parametri, sarebbe stata una soluzione semplice, però non avrebbe permesso l'utilizzo di metodi come isAlive e getDeathReason
\end{itemize}



\subsubsection{Design di ReportImpl come componente JavaFX riutilizzabile}

\textbf{Problema:} Alla fine di ogni semestre e alla fine della partita si vuole presentare un report al giocatore nel quale sono presenti i valori di ogni parametro e delle frasi di storytelling.

\noindent\textbf{Soluzione:} \texttt{ReportImpl} è un componente della View e implementa l'interfaccia \texttt{Report}. È progettato per essere riutilizzato sia come schermata di riepilogo inter-semestrale sia come schermata di fine partita, senza creare istanze multiple.

\begin{figure}[H]
    \centering
    \fbox{\includegraphics[width=0.8\textwidth]{UML-Piraino-ReportGUI.png}}
    \caption{Design di ReportImpl come componente riutilizzabile}
    \label{fig:report}
\end{figure}

\noindent Alcune alternative sono:
\begin{itemize}
    \item Pattern Strategy: Definire un'interfaccia comune per il report e realizzare implementazioni separate per le diverse tipologie di riepilogo. Questa alternativa non è stata scelta poiché le differenze grafiche e logiche tra i report sono attualmente minime, e l'introduzione del pattern avrebbe comportato solo una maggiore complessità del codice.
    \item creare elementi al livello sopra: invece di cancellare con .children().clear() si aggiunge un nuovo report sopra. Permetterebbe di creare delle animazioni in cui si mostra la modifica dei parametri, però richiederebbe molto più utilizzo di memoria.
\end{itemize}


% =============================================================================
\chapter{Sviluppo}
% =============================================================================

% -----------------------------------------------------------------------------
\section{Testing automatizzato}
% -----------------------------------------------------------------------------
Il testing automatizzato \`e stato realizzato con \textbf{JUnit~5 (Jupiter)},
con test posti in \texttt{src/test/java} e package paralleli a quelli dei
sorgenti. I test non richiedono alcuna interazione utente e sono eseguiti
automaticamente dal task \texttt{gradle build}.

I test automatici sono stati organizzati per area funzionale, in modo da coprire
le principali responsabilità dell'applicazione: la logica di avanzamento della
partita, il caricamento delle carte da configurazione esterna e l'interazione tra
Model, Controller e View.

\paragraph{Testing delle transizioni di stato del Model}
La strategia adottata \`e il \emph{State Transition Testing}: ogni test
verifica che il sistema si trovi nello stato corretto prima e dopo
un'operazione che ne altera il comportamento. L'annotazione
\texttt{@BeforeEach} garantisce l'isolamento fra test: ogni metodo riceve
un'istanza fresca del sistema sotto test, eliminando dipendenze di ordine.

Le classi di test coprono tre classi di implementazione:
\texttt{GameEngineTest} (transizioni di stato),
\texttt{GameClockTest} (progressione temporale),
\texttt{GameConfigTest} (validit\`a della configurazione).

\subparagraph{\texttt{GameEngineTest}}
\texttt{setUp()} inizializza il motore con la configurazione standard e un
mazzo vuoto:
\begin{lstlisting}
this.engine = new GameEngineImpl(GameConfigImpl.createStandard(), new Deck());
\end{lstlisting}
\begin{itemize}
    \item \texttt{testGameInitialization()}: verifica che lo stato iniziale
        sia \texttt{RUNNING} prima e dopo \texttt{engine.start()}, documentando
        che l'inizializzazione non altera lo stato logico del gioco.
    \item \texttt{testGameLostOnParameterDeath()}: simula la morte del
        parametro \texttt{FINANCES} azzerandone il livello; verifica che
        \texttt{areAllParametersAlive()} rilevi il parametro morto e che
        \texttt{getGameState()} risponda con \texttt{LOST}.
    \item \texttt{testGameWonOnTimeFinished()}: verifica la transizione
        \texttt{RUNNING} -> \texttt{WON} esaurendo il tempo tramite un ciclo
        \texttt{while (!isTimeFinished()) nextTurn();}, controllando anche
        la collaborazione fra \texttt{GameEngineImpl} e \texttt{GameClockImpl}.
\end{itemize}

\subparagraph{\texttt{GameClockTest}}
\texttt{configuration()} inizializza orologio e configurazione con
\texttt{createStandard()}.
\begin{itemize}
    \item \texttt{testSingleTurnProgression()}: verifica che una singola
        \texttt{nextTurn()} incrementi \texttt{currentTurn} di 1 senza
        alterare \texttt{currentSemester}. La cattura dei valori iniziali
        rende l'asserzione indipendente dalla configurazione specifica.
    \item \texttt{testSemesterFlow()}: verifica che dopo
        \texttt{cardsPerSemester} chiamate a \texttt{nextTurn()} il semestre
        avanzi di 1 e il contatore dei turni venga azzerato a 0.
    \item \texttt{testEndofTime()}: esegue il ciclo completo dei turni e
        verifica che \texttt{isTimeFinished()} diventi \texttt{true} al
        termine di tutti i semestri.
\end{itemize}

\subparagraph{\texttt{GameConfigTest}}
\texttt{testStandardConfigurationValidity()} verifica che
\texttt{getCardsPerSemester()} e \texttt{getSemestersPerGame()} restituiscano
valori strettamente positivi per la configurazione standard. Il contratto
coperto \`e minimo per design: verificare i valori esatti (6~e~6) renderebbe
il test fragile rispetto a future variazioni della configurazione standard.

\paragraph{Testing del caricamento delle carte da file YAML} 
I test riguardano principalmente il sistema di carte e il caricamento dei dati da file YAML. Vista la brevità ho raggruppato in un'unica classe più test.

Sono stati sottoposti a test i seguenti aspetti:
\begin{itemize}
    \item \texttt{shouldLoadCards()} e \texttt{shouldLoadFollowUp()}: verificano che le informazioni scritte nei file 
    YAML coincidano rispettivamente con ciò che è contenuto dentro l'implementazione delle carte create e dei followup.
    \item \texttt{shouldHaveUniqueIdCards()} e \texttt{shouldHaveUniqueIdChildCards()}: verificano che a gruppi le carte base e quelle figlie abbiano tutte un identificatore (\texttt{String id}) diverso.
    \item \texttt{shouldNotBeSharedIds()}: verifica che non ci siano identificatori comuni tra le carte base e quelle figlie.
    \item \texttt{shouldCharacterImageExist()}: verifica che ogni personaggio all'interno delle carte abbia il riferimento ad un'immagine da mostrare durante il gioco.
\end{itemize}

\paragraph{Testing di parametri, Controller e Observer}

La strategia adottata è la combinazione di \emph{State Transition Testing}
e \emph{Boundary Testing}: ogni test verifica il comportamento del sistema
ai confini del dominio (livelli 0 e 100 dei parametri) e nelle transizioni
di stato (vivo $\to$ morto, notifica $\to$ silenziosa dopo la morte).
L'annotazione \texttt{@BeforeEach} garantisce l'isolamento fra test,
ricreando ogni volta un'istanza fresca del sistema sotto test.
Le classi di test coprono cinque aree, per un totale di 44 test automatici:

\subparagraph{\texttt{ParameterImplTest} -- 18 test:}
Verifica il comportamento di \texttt{ParameterImpl} in tutti gli scenari
rilevanti: livello iniziale a 50, incremento e decremento del livello,
clamping agli estremi (0 e 100), transizione alive $\to$ dead, guard
clause che ignora \texttt{modify()} dopo la morte, correttezza di
\texttt{getDeathReason()} ai due estremi, fail-fast su observer nullo,
notifica dell'observer con tipo e livello corretti, assenza di
notifiche dopo la morte del parametro, e correttezza ai boundary esatti
(delta di esattamente 50 e $-$50).

\subparagraph{\texttt{ParameterTypeTest} -- 7 test:}
Verifica il pattern \emph{Type-Safe Enum con comportamento}:
tutti e quattro i valori espongono un \texttt{displayName} non nullo
e non vuoto; i nomi attesi (\texttt{"Finances"}, \texttt{"Students"},
\texttt{"Professors"}, \texttt{"Reputation"}) corrispondono ai valori
restituiti da \texttt{getDisplayName()}; il numero totale di costanti
è esattamente quattro.

\subparagraph{\texttt{ParameterObserverTest} -- 3 test:}
Verifica che \texttt{ParameterObserver} sia una vera
\texttt{@FunctionalInterface}: può essere istanziata come lambda,
riceve il tipo corretto (\texttt{ParameterType}) e il livello corretto
al momento della notifica.

\subparagraph{\texttt{PlayerInfoTest} -- 7 test:}
Verifica il pattern \emph{Value Object}: i getter restituiscono i valori
passati alla costruzione; il costruttore compatto rigetta immediatamente
\texttt{null} su ciascuno dei tre campi (fail-fast); due record costruiti
con gli stessi valori risultano uguali per \texttt{equals()}; due record
con valori diversi non lo sono; \texttt{toString()} contiene i valori
significativi.

\subparagraph{\texttt{GameControllerImplTest} -- 9 test:}
Test di integrazione che verifica la comunicazione MVC senza lanciare
JavaFX. Una \texttt{FakeView} (implementazione \emph{spy} di
\texttt{GameView}) registra le chiamate ricevute dal Controller e
permette di asserire il comportamento osservabile. I test coprono:
avvio del gioco e prima visualizzazione della carta, aggiornamento della
View dopo una decisione, notifica dell'Observer sui parametri,
robustezza a sequenze rapide di decisioni, correttezza di
\texttt{getPlayerInfo()}, range valido dei livelli restituiti da
\texttt{getCurrentParametersLevels()}, valori non nulli di
\texttt{previewDecisionDeltas()} in entrambe le direzioni, e assenza
di eccezioni prima di \texttt{startGame()}.

\subparagraph{Test della View con TestFX -- 3 classi:}
A complemento dei test di logica, la View è stata verificata con
\textbf{TestFX} e \textbf{Mockito}: \texttt{GameViewJavaFXImplTest}
(usa un mock del \texttt{GameController} per verificare che
aggiornamento parametri, tempo e schermate di fine partita non
generino eccezioni sul thread JavaFX), \texttt{ParameterIconViewTest}
(verifica la comparsa e scomparsa del pallino di anteprima al variare
di \texttt{setAffected()}) e \texttt{LoginSceneTest} (verifica login
valido e validazione dei campi vuoti, interagendo direttamente con
il grafo della scena tramite \texttt{robot.interact()}).

% -----------------------------------------------------------------------------
\section{Note di sviluppo}
% -----------------------------------------------------------------------------

% --- 3.2.1 ---------------------------------------------------------
\subsection{Pietro Dapporto}
\begin{itemize}
    \item \textbf{Libreria esterna Jackson YAML} \\
    \url{https://github.com/Federicolocaa/OOP25-Aurea/blob/dd18617aeefb0d0e4c01236c24cc8c7b8c0a804d/src/main/java/it/unibo/aurea/model/Deck.java#L31-L82}

    \item \textbf{Record Java come Data Transfer Object} - usati per \texttt{CardDTO}, \texttt{EffectDTO}, \texttt{DecisionDTO}, \texttt{FollowUpDTO}, \texttt{FollowUpsFile}, \texttt{CardsFile} di seguito solo un esempio.\\
    \url{https://github.com/Federicolocaa/OOP25-Aurea/blob/dd18617aeefb0d0e4c01236c24cc8c7b8c0a804d/src/main/java/it/unibo/aurea/model/dto/CardDTO.java#L5-L21}

    \item \textbf{Stream API e method reference} - usati anche nei test automatizzati (\texttt{CardLoaderTest}) \\
    \url{https://github.com/Federicolocaa/OOP25-Aurea/blob/dd18617aeefb0d0e4c01236c24cc8c7b8c0a804d/src/main/java/it/unibo/aurea/model/Deck.java#L58-L78}
\end{itemize}

% --- 3.2.2 ---------------------------------------------------------
\subsection{Mario Conventi}
Di seguito le feature avanzate del linguaggio Java utilizzate nella mia parte di progetto, ciascuna corredata del permalink GitHub al punto del codice in cui è impiegata:

\begin{itemize}
    \item \textbf{Libreria esterna Jackson per parsing JSON/YAML in Record} -- la configurazione del bilanciamento del gioco viene letta dal file esterno tramite \texttt{ObjectMapper} e deserializzata automaticamente in un albero di Java Record immodificabili. \\
    \url{https://github.com/Federicolocaa/OOP25-Aurea/blob/f25267cad6c0293387dbc32852e679f9b4ac103c/src/main/java/it/unibo/aurea/model/DifficultySettings.java#L45-L51}

    \item \textbf{Gestione sicura dei null tramite \texttt{Optional}} -- l'estrazione delle carte figlie (\emph{FollowUp}) forzate dalla coda e la gestione dei \texttt{findFirst()} negli stream restituiscono e manipolano istanze di \texttt{Optional<Card>}, evitando in modo robusto eccezioni di tipo \texttt{NullPointerException}. \\
    \url{https://github.com/Federicolocaa/OOP25-Aurea/blob/f25267cad6c0293387dbc32852e679f9b4ac103c/src/main/java/it/unibo/aurea/model/FollowUpQueue.java#L38-L62}

    \item \textbf{Stream API per pipeline di filtraggio} -- le code e i selettori di carte applicano filtri multipli e metodi come \texttt{allMatch} o \texttt{forEach} su \texttt{Stream} per isolare e registrare eventi in logiche compatte e dichiarative. \\
    \url{https://github.com/Federicolocaa/OOP25-Aurea/blob/f25267cad6c0293387dbc32852e679f9b4ac103c/src/main/java/it/unibo/aurea/model/FollowUpQueue.java#L74-L77}

    \item \textbf{Copia difensiva (Defensive Copy) per incapsulamento} -- l'orchestratore principale del gioco nasconde il proprio stato interno fornendo unicamente copie immodificabili tramite \texttt{List.copyOf()}, proteggendo le collezioni e garantendo l'immutabilità richiesta esternamente. \\
    \url{https://github.com/Federicolocaa/OOP25-Aurea/blob/f25267cad6c0293387dbc32852e679f9b4ac103c/src/main/java/it/unibo/aurea/model/GameEngineImpl.java#L114-L116}

    \item \textbf{Static Nested Class} -- il concetto di evento in attesa temporizzato all'interno della coda delle conseguenze è racchiuso privatamente all'interno di una classe annidata statica (\texttt{ActiveFollowUp}), oscurando i dettagli implementativi all'esterno. \\
    \url{https://github.com/Federicolocaa/OOP25-Aurea/blob/f25267cad6c0293387dbc32852e679f9b4ac103c/src/main/java/it/unibo/aurea/model/FollowUpQueue.java#L80-L100}
\end{itemize}

\paragraph{Codice di terze parti}
Nessuna porzione di codice della mia parte di progetto è stata tratta da fonti esterne. Tra le librerie impiegate, si segnala l'uso di Jackson per la lettura dei file di configurazione YAML.

% --- 3.2.3 ---------------------------------------------------------
\subsection{Federico Locatelli}
Di seguito le feature avanzate del linguaggio e dell'ecosistema Java
utilizzate nella mia parte di progetto, ciascuna corredata del permalink
GitHub al punto del codice in cui è impiegata.

\begin{itemize}
    \item \textbf{\texttt{@FunctionalInterface}} -- \texttt{ParameterObserver}
        è dichiarata come interfaccia funzionale: può essere istanziata
        come lambda o method reference senza classi anonime. \\
        \url{https://github.com/Federicolocaa/OOP25-Aurea/blob/f25267cad6c0293387dbc32852e679f9b4ac103c/src/main/java/it/unibo/aurea/model/api/ParameterObserver.java#L8-L16}

    \item \textbf{Method reference come observer} -- la View viene collegata
        al Model tramite \texttt{this.view::updateSingleParameter}, passata
        come \texttt{ParameterObserver} senza lambda esplicita. \\
        \url{https://github.com/Federicolocaa/OOP25-Aurea/blob/f25267cad6c0293387dbc32852e679f9b4ac103c/src/main/java/it/unibo/aurea/controller/GameControllerImpl.java#L48-L50}

    \item \textbf{Interfacce funzionali \texttt{BiConsumer} e \texttt{Function}} --
        i tre comportamenti variabili di \texttt{CardPanel} sono rappresentati
        da interfacce funzionali generiche iniettate dall'esterno, disaccoppiando il
        componente visuale dal Controller. \\
        \url{https://github.com/Federicolocaa/OOP25-Aurea/blob/f25267cad6c0293387dbc32852e679f9b4ac103c/src/main/java/it/unibo/aurea/view/CardPanel.java#L92-L93}
        
    \item \textbf{\texttt{TextFormatter} con lambda} -- i campi di testo
        del login limitano la lunghezza dell'input tramite un predicato
        lambda passato al \texttt{TextFormatter} di JavaFX. \\
        \url{https://github.com/Federicolocaa/OOP25-Aurea/blob/f25267cad6c0293387dbc32852e679f9b4ac103c/src/main/java/it/unibo/aurea/view/LoginScene.java#L85-L93}

    \item \textbf{Switch expression Java~21} -- la selezione della difficoltà
        converte la stringa della \texttt{ChoiceBox} nell'enum
        \texttt{Difficulty} tramite switch expression con arrow syntax. \\
        \url{https://github.com/Federicolocaa/OOP25-Aurea/blob/f25267cad6c0293387dbc32852e679f9b4ac103c/src/main/java/it/unibo/aurea/view/LoginScene.java#L116}

    \item \textbf{JavaFX \texttt{Timeline} / \texttt{KeyFrame} / \texttt{KeyValue}} --
        l'animazione di riempimento dell'icona parametro interpola
        fluidamente il livello tramite \texttt{Interpolator.EASE\_BOTH}. \\
        \url{https://github.com/Federicolocaa/OOP25-Aurea/blob/f25267cad6c0293387dbc32852e679f9b4ac103c/src/main/java/it/unibo/aurea/view/ParameterIconView.java#L150-L152}

    \item \textbf{JavaFX \texttt{ParallelTransition}} -- l'animazione di
        ingresso della carta compone in parallelo una
        \texttt{TranslateTransition} (scivolamento verticale) e una
        \texttt{FadeTransition} (dissolvenza), ottenendo un effetto
        fluido con una sola chiamata a \texttt{play()}. \\
        \url{https://github.com/Federicolocaa/OOP25-Aurea/blob/f25267cad6c0293387dbc32852e679f9b4ac103c/src/main/java/it/unibo/aurea/view/CardPanel.java#L342}

    \item \textbf{JavaFX \texttt{BlendMode.SCREEN}} -- le icone dei parametri
        usano \texttt{BlendMode.SCREEN} per sovrapporre il layer attivo
        a quello dimmed senza mascherare il canale alpha dell'immagine. \\
        \url{https://github.com/Federicolocaa/OOP25-Aurea/blob/f25267cad6c0293387dbc32852e679f9b4ac103c/src/main/java/it/unibo/aurea/view/ParameterIconView.java#L87-L90}

    \item \textbf{\texttt{Platform.runLater}} -- tutti i metodi della View
        che modificano nodi della scena delegano l'esecuzione al thread
        JavaFX, prevenendo eccezioni di accesso concorrente. \\
        \url{https://github.com/Federicolocaa/OOP25-Aurea/blob/f25267cad6c0293387dbc32852e679f9b4ac103c/src/main/java/it/unibo/aurea/view/GameViewJavaFXImpl.java#L322-L333}

    \item \textbf{\texttt{SVGPath} per grafica vettoriale} -- l'icona di
        chiusura è costruita tramite \texttt{SVGPath} con path data
        inline, scalabile a qualsiasi risoluzione senza immagini esterne. \\
        \url{https://github.com/Federicolocaa/OOP25-Aurea/blob/f25267cad6c0293387dbc32852e679f9b4ac103c/src/main/java/it/unibo/aurea/view/GameViewJavaFXImpl.java#L193}
    \item \textbf{JavaFX \texttt{MediaPlayer} e \texttt{AudioClip}} -- 
        il sistema audio usa \texttt{MediaPlayer} per la musica di 
        sottofondo in loop e \texttt{AudioClip} per gli effetti brevi 
        a bassa latenza; \texttt{setOnEndOfMedia} riprende il sottofondo 
        al termine di un suono speciale. \\
        \url{https://github.com/Federicolocaa/OOP25-Aurea/blob/f25267cad6c0293387dbc32852e679f9b4ac103c/src/main/java/it/unibo/aurea/view/AudioManager.java#L109-L119}
\end{itemize}

\paragraph{Codice di terze parti}
Nessuna porzione di codice della mia parte di progetto è stata tratta
da fonti esterne. Le librerie utilizzate, oltre a Java~21 e JUnit~5,
sono JavaFX per l'intero layer View e SpotBugs per l'annotazione
\texttt{@SuppressFBWarnings} nel Controller.

% --- 3.2.4 ---------------------------------------------------------
\subsection{Nikita Piraino}
Di seguito le feature avanzate del linguaggio e dell'ecosistema Java utilizzate
nella mia parte di progetto, ciascuna corredata del permalink GitHub al punto
del codice in cui \`e impiegata.
\begin{itemize}
    \item \textbf{Utilizzo di generici bounded}\\
    Permalink1: \url{https://github.com/Federicolocaa/OOP25-Aurea/blob/73489c086c7d02bec683ad10117f1d45959b317f/src/main/java/it/unibo/aurea/model/CardSelector.java#L30} \\
    permalink2: \url{https://github.com/Federicolocaa/OOP25-Aurea/blob/73489c086c7d02bec683ad10117f1d45959b317f/src/main/java/it/unibo/aurea/model/CardSelector.java#L87} \\
    permalink3: \url{https://github.com/Federicolocaa/OOP25-Aurea/blob/73489c086c7d02bec683ad10117f1d45959b317f/src/main/java/it/unibo/aurea/model/CardSelector.java#L92} \\
    permalink4: \url{https://github.com/Federicolocaa/OOP25-Aurea/blob/73489c086c7d02bec683ad10117f1d45959b317f/src/main/java/it/unibo/aurea/model/CardSelector.java#L107}
    \item \textbf{Lambda expressions} \\
    Permalink1: \url{https://github.com/Federicolocaa/OOP25-Aurea/blob/73489c086c7d02bec683ad10117f1d45959b317f/src/main/java/it/unibo/aurea/view/ReportImpl.java#L64} \\
    Permalink2: \url{https://github.com/Federicolocaa/OOP25-Aurea/blob/73489c086c7d02bec683ad10117f1d45959b317f/src/main/java/it/unibo/aurea/view/InfoButtonImpl.java#L39} \\
    Permalink3: \url{https://github.com/Federicolocaa/OOP25-Aurea/blob/73489c086c7d02bec683ad10117f1d45959b317f/src/main/java/it/unibo/aurea/view/InfoButtonImpl.java#L64}  \\
    \item \textbf{Method Reference} \\
    permalink1: \url{https://github.com/Federicolocaa/OOP25-Aurea/blob/9d158606c6f5d3f9b297de450441bb3df5c1d139/src/test/java/it/unibo/aurea/model/GameClockTest.java#L71}
    \item \textbf{Stream} \\
    Permalink1: \url{https://github.com/Federicolocaa/OOP25-Aurea/blob/73489c086c7d02bec683ad10117f1d45959b317f/src/main/java/it/unibo/aurea/model/GameEngineImpl.java#L126} \\
    \item \textbf{Libreria di terze parti - JavaFX} \\
    i seguenti permalink puntano a delle righe qualsiasi, perchè le seguenti classi sono state realizzate interamente da me. Alcuni metodi sono stati copiati da codice che Federico Locatelli (compagno) 
    aveva scritto in altre classi sostituite da quelle seguenti.
    permalink1: \url{https://github.com/Federicolocaa/OOP25-Aurea/blob/73489c086c7d02bec683ad10117f1d45959b317f/src/main/java/it/unibo/aurea/view/EndgameOverlay.java#L48} \\
    permalink2: \url{https://github.com/Federicolocaa/OOP25-Aurea/blob/73489c086c7d02bec683ad10117f1d45959b317f/src/main/java/it/unibo/aurea/view/ReportImpl.java#L19} \\
    permalink3: \url{https://github.com/Federicolocaa/OOP25-Aurea/blob/73489c086c7d02bec683ad10117f1d45959b317f/src/main/java/it/unibo/aurea/view/InfoButtonImpl.java#L20}
\end{itemize}

\paragraph{Codice di terze parti}
Nessuna porzione di codice della mia parte di progetto \`e stata tratta da
fonti esterne.

% =============================================================================
\chapter{Commenti finali}
% =============================================================================

% -----------------------------------------------------------------------------
\section{Autovalutazione e lavori futuri}
% -----------------------------------------------------------------------------

% --- Pietro -----------------------------------------------------
\subsection{Pietro Dapporto}
\paragraph{Contributi principali:}
Il mio ruolo all'interno del gruppo ha riguardato soprattutto la struttura di base del gioco 
e la modellazione delle carte. Per questo motivo, in alcune parti del mio contributo non è 
stato necessario utilizzare costrutti Java particolarmente complessi: l'obiettivo principale 
era costruire una struttura chiara, coerente e facilmente estendibile.

\paragraph{Aspetti di cui sono soddisfatto:}
La parte più creativa del lavoro, di cui sono particolarmente soddisfatto, è stata la 
creazione delle carte di gioco. Dal punto di vista implementativo, questa attività non ha 
richiesto tecnicismi avanzati, ma ha comportato un lavoro importante di progettazione del 
contenuto: definire le situazioni proposte al giocatore, le possibili risposte e i pesi 
associati agli effetti sui parametri del gioco. Ritengo inoltre di aver adottato soluzioni 
strutturali discrete, soprattutto per quanto riguarda la separazione tra i dati delle carte 
in file di configurazione esterni e la logica principale dell'applicazione.

Considero positivo il fatto di aver progressivamente migliorato l'architettura del sistema, 
passando da soluzioni inizialmente più semplici ma meno scalabili a una struttura più modulare 
e mantenibile. 

\paragraph{Aspetti migliorabili e lavori futuri:}
Per quanto riguarda i punti deboli, il file YAML contenente le carte è diventato piuttosto 
esteso. Con l'aggiunta di altre carte la gestione di un unico file può diventare meno comoda. 
Ho valutato la possibilità di suddividerlo in più file, ma ho ritenuto che, almeno in questa 
fase del progetto, mantenere un unico file fosse più semplice da gestire e riducesse il 
rischio di frammentare eccessivamente le informazioni.

Un altro limite riguarda il livello di astrazione introdotto. In alcuni punti la struttura 
del codice è stata pensata per supportare possibili estensioni future, come modelli diversi 
di carte o di effetti, che però non sono state effettivamente sviluppate nel progetto finale. 
Lo stesso vale per l'utilizzo dei DTO, che separano bene la rappresentazione esterna dei dati 
dal model, ma introducono anche un passaggio intermedio aggiuntivo.

% --- Mario ------------------------------------------------------
\subsection{Mario Conventi}

\paragraph{Contributi principali:}
Il mio ruolo all'interno del progetto si è concentrato sulla progettazione e sullo sviluppo del nucleo computazionale del Model. Mi sono occupato della realizzazione del \texttt{GameEngineImpl}, l'orchestratore principale del gioco, e dei suoi sotto-componenti algoritmici specializzati: la selezione adattiva delle carte (\texttt{CardSelector}), il sistema asincrono per l'accodamento delle carte figlie (\texttt{FollowUpQueue}), e la gestione dinamica del bilanciamento tramite le impostazioni di difficoltà (\texttt{DifficultySettings}). Il mio obiettivo primario è stato trasformare una semplice macchina a stati in un motore di gioco adattivo e configurabile, in grado di rispondere in modo bilanciato alle scelte del giocatore.

\paragraph{Aspetti di cui sono soddisfatto:}
Sono particolarmente orgoglioso delle soluzioni adottate per risolvere i problemi architetturali legati agli algoritmi principali. L'implementazione del \texttt{CardSelector} con un sistema di selezione pesata dinamica garantisce un'ottima fluidità di gioco, prevenendo la frustrazione senza sfociare in un determinismo assoluto. Ritengo particolarmente riuscita la logica matematica alla base del bilanciamento: aver parametrizzato dinamicamente i pesi di estrazione e i moltiplicatori d'impatto tramite \texttt{DifficultySettings} permette al gioco di adattare la curva di difficoltà in modo elegante, modulando il "soccorso" algoritmico offerto al giocatore nelle situazioni critiche senza mai dover alterare il codice \textit{core} del motore.

\paragraph{Aspetti migliorabili e lavori futuri:}
Rivedendo il lavoro svolto, sebbene abbia cercato di applicare rigorosamente il Single Responsibility Principle (SRP) estraendo svariati sotto-componenti, il \texttt{GameEngineImpl} mantiene ancora una forte centralità come coordinatore; un possibile sviluppo futuro potrebbe prevedere l'introduzione di pattern più avanzati (come Event Bus o Mediator) per sgravare ulteriormente l'engine dalle interazioni troppo fitte. Inoltre, come sviluppo futuro per l'algoritmo di selezione delle carte, sarebbe interessante introdurre un sistema di "memoria a breve termine" per impedire che eventi tematicamente simili (ad esempio due eventi consecutivi legati ai professori) vengano estratti a distanza troppo ravvicinata, migliorando ancor di più la narrazione emergente del gioco.

% --- Federico ---------------------------------------------------
\subsection{Federico Locatelli}

\paragraph{Il mio ruolo nel gruppo:}
Fin dall'inizio mi sono ritrovato a gravitare verso la parte
dell'applicazione che si vede e si tocca: la View. È la dimensione del
progetto che mi diverte di più, perché mi costringe a ragionare non solo
su ``cosa fa il programma'' ma su come l'esperienza arriva a chi gioca ---
quanto è fluida un'animazione, se un suono cade nel momento giusto, se
un'icona comunica subito ciò che deve. Mi sono occupato dell'intero layer
grafico, delle astrazioni reattive del Model che lo alimentano
(\texttt{Parameter}, \texttt{ParameterObserver}, \texttt{ParameterImpl})
e dell'estensione del Controller con i metodi di anteprima. Ho infine
curato buona parte del testing automatizzato. Tengo molto alla cura del
dettaglio, e questa parte del progetto mi ha permesso di esprimerla.

\paragraph{Aspetti di cui sono soddisfatto:}
La parte di cui sono più soddisfatto è il sistema reattivo dei parametri. Aver
reso \texttt{ParameterObserver} una \texttt{@FunctionalInterface} e averla
collegata tramite method reference nel Controller mi ha permesso di ottenere un
collegamento Model--View pulito, in cui ogni componente conosce solo ciò che gli
serve. Ritengo riuscito anche il disaccoppiamento della carta interattiva dalla
logica di gioco: il componente visuale non ha alcun riferimento al Controller,
eppure ne pilota il comportamento attraverso comportamenti iniettati
dall'esterno. Sul piano della presentazione, considero positiva la coerenza
visiva e sonora dell'applicazione: il foglio di stile unifica il tema
dell'Ateneo in tutti i componenti, mentre il sistema audio ha dato al gioco un
carattere che nelle prime versioni mancava.

\paragraph{Aspetti migliorabili e lavori futuri:}
Un primo limite riguarda la gestione del threading JavaFX: attualmente ogni
metodo della View invoca \texttt{Platform.runLater()} per conto proprio, e a
posteriori avrei centralizzato questa responsabilità in un wrapper dedicato per
ridurre la ripetizione. Un secondo limite è il foglio di stile, cresciuto in
modo organico man mano che aggiungevo componenti: una suddivisione in moduli
tematici, pianificata fin dall'inizio, lo avrebbe reso più mantenibile. Sono
scelte di struttura che, ho imparato, conviene affrontare presto, prima che il
codice cresca. Tra i possibili sviluppi futuri mi interesserebbe introdurre
animazioni sui valori numerici dei parametri durante l'aggiornamento e
completare l'integrazione della difficoltà nella View con un feedback visivo
dedicato alla modalità \emph{Hard}.

% --- Nikita -----------------------------------------------------
\subsection{Nikita Piraino}
Sono un ragazzo a cui piace molto l'azione di ``unire i puntini'' e il ruolo che ho scelto mi ha permesso di mettere ulteriormente in pratica le conoscenze necessarie per farlo bene. Ho avuto la possibilità di ripassare quasi ogni aspetto della materia studiata. \\

Sono \textbf{soddisfatto} del lavoro che ho svolto perchè nella scrittura dei miei codici ho identificato velocemente una sequenza da seguire per evitare di  fare  modifiche ``scomode'' in futuro. Quando ho toccato  codici per la maggior parte scritti dai miei compagni non mi sono trovato in difficoltà nel comprenderne il funzionamento(ignorando ovviamente le lacune riguardanti la non conoscenza iniziale delle liberie). \\
Un aspetto per cui mi è stato molto utile lavorare a questo progetto è stato quello della documentazione: mi piace come attività, mi sono reso conto di quanto scriverla mi fornisse spunti per migliorare il codice; quindi \textbf{mi pento} di averla fatta per la maggior parte  verso la fine in quanto mi ha portato a voler rivedere alcune scelte e quindi riscriverla. La mia responsabilità principale era creare del codice che fosse chiaro, estendibile per nuove modifiche, adatto per le necessità di ogni mio compagno e futuro sviluppatore, quindi non ho avuto la necessità di creare costrutti particolarmente articolati dal punto di vista di \textit{feature avanzate di linguaggio}.\\

Sempre legato all'unire i puntini sento di \textbf{essere stato} molto un connettore all'interno del gruppo, chiedevo aggiornamenti, organizzavo le riunioni che facevamo per sintonizzarci e se scoprivo qualcosa di interessante lo condividevo anche agli altri.

% =============================================================================
\appendix
% =============================================================================

% =============================================================================
\chapter{Guida utente}
% =============================================================================
Questa guida accompagna l'utente nei passaggi essenziali per avviare il
programma e giocare una partita ad \textit{Aurea Mediocritas}.

\section{Avvio dell'applicazione}

L'applicazione può essere avviata in due modi, a partire dalla cartella radice
del progetto. Il primo, pensato per chi ha a disposizione l'ambiente di
sviluppo, esegue direttamente il gioco tramite Gradle:
\begin{lstlisting}[language=bash]
./gradlew run
\end{lstlisting}
In alternativa è possibile generare un archivio eseguibile autocontenuto (che
include JavaFX e tutte le dipendenze) e lanciarlo con Java:
\begin{lstlisting}[language=bash]
./gradlew shadowJar
java -jar build/libs/Aurea-all.jar
\end{lstlisting}
È richiesto un Java Development Kit versione 21 o successiva.

\section{Schermata iniziale e creazione della partita}

All'avvio compare la schermata iniziale. Prima di cominciare, l'utente deve
inserire il nome del Rettore e quello della facoltà (o università), e
selezionare il livello di difficoltà desiderato fra quelli proposti. Il
pulsante delle informazioni, in basso a destra, apre una breve spiegazione
delle regole. Una volta compilati i campi, la partita si avvia premendo il
pulsante \textit{Begin your Reign}.

\begin{figure}[H]
    \centering
    \includegraphics[width=0.55\textwidth]{schermata_iniziale.png}
    \caption{Schermata iniziale: inserimento dei dati, scelta della difficoltà e avvio della partita.}
\end{figure}

\section{Schermata di gioco}

A partita avviata, al centro dello schermo compare una carta: ciascuna
rappresenta una situazione a cui il Rettore deve rispondere con una decisione.
Nella parte superiore sono mostrati i quattro parametri dell'Ateneo
(Finanze, Studenti, Docenti e Reputazione), ognuno con un valore compreso fra
0 e 100. L'obiettivo è mantenerli tutti in equilibrio fino al termine del
mandato: se anche uno solo di essi raggiunge il minimo o il massimo, la partita
si conclude con una sconfitta.

\begin{figure}[H]
    \centering
    \includegraphics[width=0.60\textwidth]{schermata_gioco.png}
    \caption{Schermata principale di gioco, con la carta corrente e i parametri dell'Ateneo.}
\end{figure}

\section{Scelta delle decisioni}

Per rispondere a una carta, l'utente la tiene premuta e la trascina verso
destra o verso sinistra: il trascinamento rivela progressivamente le due
risposte possibili. Durante il movimento, alcuni indicatori compaiono sopra i
parametri per anticipare quali verranno influenzati dalla scelta in corso. Al
rilascio della carta la decisione viene confermata, i parametri si aggiornano e
il gioco propone la carta successiva.

\begin{figure}[H]
    \centering
    \includegraphics[width=0.65\textwidth]{swipe_destra.png}
    \caption{Swipe a destra: approvazione della proposta.}
\end{figure}

\begin{figure}[H]
    \centering
    \includegraphics[width=0.65\textwidth]{swipe_sinistra.png}
    \caption{Swipe a sinistra: rifiuto della proposta.}
\end{figure}

\section{Riepilogo e fine della partita}

Al termine di ogni semestre viene mostrata una schermata di riepilogo con i
valori correnti dei quattro parametri, utile per valutare l'andamento del
mandato prima di proseguire.

\begin{figure}[H]
    \centering
    \includegraphics[width=0.70\textwidth]{riepilogo_semestre.png}
    \caption{Schermata di riepilogo mostrata al termine di ogni semestre.}
\end{figure}

La partita termina in due casi: quando il Rettore completa l'intero mandato
(vittoria) oppure quando uno dei parametri raggiunge un valore critico
(sconfitta). In entrambi i casi compare una schermata finale con l'esito,
dalla quale è possibile iniziare una nuova partita senza riavviare
l'applicazione.

\begin{figure}[H]
    \centering
    \includegraphics[width=0.70\textwidth]{schermata_vittoria.png}
    \caption{Schermata di vittoria: il Rettore ha completato il mandato.}
\end{figure}

\begin{figure}[H]
    \centering
    \includegraphics[width=0.70\textwidth]{schermata_sconfitta.png}
    \caption{Schermata di sconfitta: un parametro ha raggiunto un valore critico.}
\end{figure}

% =============================================================================
\chapter{Esercitazioni di laboratorio}
% =============================================================================

\section*{pietro.dapporto@studio.unibo.it}
\begin{itemize}
    \item Lab 09: \url{https://virtuale.unibo.it/mod/forum/discuss.php?d=208718#p287289}
    \item Lab 11: \url{https://virtuale.unibo.it/mod/forum/discuss.php?d=208718#p287289}
\end{itemize}

\section*{nikita.piraino@studio.unibo.it}
\begin{itemize}
    \item Lab 06: \url{https://virtuale.unibo.it/mod/forum/discuss.php?d=206731#p284135}
    \item Lab 07: \url{https://virtuale.unibo.it/mod/forum/discuss.php?d=207193#p285046}
    \item Lab 08: \url{https://virtuale.unibo.it/mod/forum/discuss.php?d=207921#p285907}
    \item Lab 09: \url{https://virtuale.unibo.it/mod/forum/discuss.php?d=208718#p287175}
    \item Lab 10: \url{https://virtuale.unibo.it/mod/forum/discuss.php?d=209589#p288241}
    \item Lab 11: \url{https://virtuale.unibo.it/mod/forum/discuss.php?d=210617#p289551}
    \item Lab 12: \url{https://virtuale.unibo.it/mod/forum/discuss.php?d=211539#p290514}
\end{itemize}

\section*{federico.locatelli3@studio.unibo.it}
\begin{itemize}
    \item Lab 8: \url{https://virtuale.unibo.it/mod/forum/discuss.php?d=207921#p286195}
    \item Lab 9: \url{https://virtuale.unibo.it/mod/forum/discuss.php?d=208718#p287053}
\end{itemize}

\end{document}